% Modernised reading — fonds Alexandre Grothendieck, shelfmark 62, pages 1-96.
%
% This is an INTERPRETATION, not a transcription. It re-reads the pages in
% today's notation and vocabulary, states the hypotheses the manuscript leaves
% implicit, and drops the critical apparatus so that the mathematics reads
% straight through. Everything the brackets and underlines carried moves into a
% footnote. It is held to being mathematically correct as it stands: where the
% page is loose or mistaken, the true statement is what appears here and the
% footnote says what the page has. In any dispute of substance the
% transcription governs: whoever cites Grothendieck must cite batch-01.fr.tex
% to batch-05.fr.tex.
%
% Pass: Opus 5.5 (claude-opus-5-5), 2026-10-03 — first pass, unchecked against
% the pages by a human. Not measured against the Opus 5 reference reading of
% folder 115. The folder was read whole — five batches, pages 1-96 — before any
% of this was written, and it is written as ONE file for the shelfmark.
%
% Coverage. Every page that carries his mathematics in the transcriptions is
% read here, with one exception: page 72. Batch 4's header announces « p. 72 a
% third recapitulation of sn, formulas (1)-(6) » and says no page was skipped,
% but its body has no \page{72} (it goes from 71 to 73). The reading cannot
% cover what the transcription does not give; this is flagged in the spine
% section and at its place. Pages 1, 2, 17, 19, 25, 26, 28, 30, 32, 34, 36,
% 38, 40, 44, 48, 50, 52, 54, 55, 58 are covers, blanks or unrelated typed
% versos, as the transcriptions record.
%
% Two moutures of one subject. Pages 1-57 are the algebraic theory over an
% arbitrary base (equations, moduli, the anharmonic group, j); pages 58-96,
% under their own cover « Théorie transcendante (Généralités classiques) »,
% are the classical analytic theory over C, itself in three layers (59-60 an
% older hand; 61-73 a first exposition; 74-96 a second, numbered 1)-5) then
% Θ, J, Legendre). The layers are kept distinct; the spine says so.
%
% NOTATION fixed for the whole folder (stated in the spine section):
%   pi : C -> S the curve (his f), e its unit section (his s), D = e(S) the
%     divisor; Lie sheaf \mathcal{L} = Lie(C/S) (his \mathcal{L} / L_{C/S}),
%     omega = \mathcal{L}^\vee; E_n = pi_* O_C(nD); gr_n E = E_n/E_{n-1}
%     (his \mathcal{L}_n);
%   x of pole order 2, y of pole order 3 (MODERN convention; pages 8-16 call
%     them y and x, page 24 and 56 already use the modern one);
%   the involution (t-2)/(2t-1) of page 27 is kappa here (his sigma, which is
%     kept for Weierstrass's sigma function); elements of S_3 in general g, h;
%     the transpositions keep his lambda, mu, nu; the middle coordinate of
%     page 35 is w (his s); the quotient X -> P^1_S of pages 41-47 is q (his
%     p and f); the S_3-torsor of page 49 is T (his S'), S' = S[sqrt(-3)] kept;
%   j is HIS invariant throughout pages 18-53, j = 1 - J = 1 - j_cl/1728 with
%     J Klein's g_2^3/Delta (page 91) and j_cl the usual invariant; so j = 0 is
%     the curve with automorphisms of order 4, j = 1 the one of order 6;
%   half-periods: omega_1 + omega_2 + omega_3 = 0 (pages 62, 90); page 78's
%     omega_3 = omega_1 + omega_2 is the same point of C/P (different sign of
%     eta_3), flagged there; P the period lattice (his gothic P).
%
% Substantive departures from the pages, each footnoted where it occurs:
%   - p. 7: L_{n+1} = L^{n+1} holds for n >= 1 (boxed « n >= 2 »);
%   - p. 6: « E_3 engendre Γ_* » read as: E_3 generates the algebra of
%     functions with poles along D, and the Veronese ring ⊕ E_{3k};
%   - p. 9: the y z^2 term missing from the homogenised cubic is restored;
%     « P^3 » read P^2; margin « Δ ∈ L^{⊗12} » is ω^{⊗12} = L^{⊗-12};
%   - p. 10: « δ ≠ 0 » read « δ invertible », as the page says just above;
%   - p. 11: BXY^2 read BXY (as p. 9 has it); p. 12: S' -> S is smooth
%     surjective of finite type, not finite;
%   - p. 13: the second « x » is y;
%   - p. 18: the cubic x^3 + 27 (j-1)/j (x+1) is not a normal form of
%     invariant j; the correct one is page 24's 4x^3 - 27 (j-1)/j (x+1);
%     p. 20's « Δ = b^3 - 27c^2 ? » holds for 4x^3 - bx - c, not x^3 + bx + c
%     (whose discriminant is -4b^3 - 27c^2);
%   - p. 22: the identity in the roots holds up to sign (irrelevant for j);
%   - p. 27: « Normalisateur : S_6 » read as page 29's S_3 x Z/2 (S_6 is not
%     a subgroup of PGL_2);
%   - p. 27: the multiplication table is right for the convention gh = « g
%     then h »; stated;
%   - p. 43: q is the quotient map, of degree 6, not an isomorphism; fibres
%     over infinity, 0, 1 are X', X'', the bisection — as pages 39, 43 (j in
%     P^1 - infinity) and 47 force, against the partly illegible assignment
%     on page 43;
%   - p. 62/67: Legendre's relation is +iπ/2 for Im(ω_2/ω_1) > 0; the labels
%     pages 65-69 force (ω_1 = iK', ω_2 = K) are the other orientation;
%   - p. 63: σ is odd, not even;
%   - p. 64: Jacobi's formula with exponent -2(n_1η_1+n_2η_2)z (sign);
%   - p. 67-69: poles of sn, cn, dn are at ±ω_1 (not ω_2), zeros of sn at 0,
%     2ω_2; cn(z+2ω_2) = -cn, cn(z+2ω_3) = cn; periods of cn 4ω_2, 2ω_3, of dn
%     2ω_2, 4ω_1; residues of cn at ω_1 is -i/k, of dn -i (checked
%     numerically); k = ± e^{η_1ω_1}/σ(ω_1)^2 (page reads σ'ω_2);
%   - p. 69: cn solves the sn-equation for ρ^2 = -k^2/(1-k^2), not
%     ρ = k/sqrt(1-k^2);
%   - p. 74: a') is the residue theorem on a period parallelogram, not
%     Liouville; p. 74 (2) defines a meromorphic, not holomorphic, function;
%   - p. 75: « 2P » read P; « le polynôme est unitaire » is false (leading
%     coefficient c_k);
%   - p. 80: ζ = 1/u - (c_1/3)u^3 - ... (page c_2/3);
%   - p. 82 (16): exponent ... + iπ (page 2iπ, which loses the sign);
%   - p. 86 (8): a factor t is missing; (10), (11), (12) are right (checked);
%     the constant of p. 87 is i sqrt(i), not sqrt(i);
%   - p. 88: e^{-π} ≈ 1/23; the reduction Re T > 1 is for real T;
%   - p. 88 (14): Θ_2 carries (-1)^n in its cosine form;
%   - p. 89 (15): e_3 = (1/3)(π/2ω_1)^2 (Θ_1(0)^4 - Θ_2(0)^4);
%   - p. 90: g_2 = -4(e_2e_3 + e_3e_1 + e_1e_2); g_3 = 28 c_2; e_i series
%     needs its -1/w^2 terms; the sums for c_1, c_2 run over non-zero periods;
%   - p. 91: g_2, g_3 are homogeneous of degree -4, -6;
%   - p. 92: J(i) = 1 (g_3 = 0), J(ρ) = 0 (g_2 = 0) — the page's values 0, 1
%     are those of 1 - J, i.e. of HIS j of pages 39-47;
%   - p. 93: dn = 1/ch for k = 1;
%   - p. 94: zeros of cn at K + 2mK + 2m'iK'; dn(u+K'') = -ik' sn/cn (checked
%     numerically; page +); margin: sn has real zeros 2mK;
%   - p. 95: 3e_2 = (k'^2 - 2)M^2.
%
% Announced and never established: the proof of Proposition 3 (p. 6-7, the
% block is cancelled); the explicit discriminant (p. 10, « Ici le
% discriminant s'explicite : ... »); the normalisation asked for at the foot
% of p. 16; « Prouver : le normalisateur ... » (p. 29); the supplementary data
% needed when j is not invertible (p. 45); the elimination « on peut aussi
% éliminer ... » (p. 53); the quotients by κ « il serait facile de ... »
% (p. 39); the value of sn ω_2 on p. 67 (denominator illegible); page 72.
\documentclass[11pt,a4paper]{article}
\input{../preamble/grothendieck}

\folder{62}
\pages{1}{96}
\dating{[à partir de 1963]}
\watermark{Édition de démonstration\\interprétation personnelle de l'œuvre}
\foldertitle{Courbes elliptiques (généralités - vieille rédaction) : notes manuscrites (s.d.) — lecture modernisée du dossier entier}

\begin{document}

\section*{Résumé}

\begin{resume}
Ces quatre-vingt-seize pages parlent des courbes elliptiques : les courbes
qu'on écrit $y^2 = x^3 + ax + b$, et qui sont en même temps des groupes --- on
sait y additionner deux points. Elles demandent comment on les classe.

Sur les nombres complexes, la réponse est ancienne et c'est la seconde moitié
du dossier qui la rappelle. Une courbe elliptique complexe est un tore : le plan
$\mathbf{C}$ divisé par un réseau de périodes. Les fonctions sur ce tore sont
les fonctions « doublement périodiques » de Weierstrass et de Jacobi, $\wp$,
$\zeta$, $\sigma$, $\mathrm{sn}$, $\mathrm{cn}$, $\mathrm{dn}$, les fonctions
thêta ; et deux tores sont isomorphes exactement quand un seul nombre, la
fonction modulaire $J$, prend la même valeur. Ces pages-là sont un formulaire
classique, rédigé deux fois et plus soigneusement la seconde.

Le problème des cinquante premières pages est de faire la même chose quand il
n'y a plus de nombres complexes, plus de réseau, plus d'intégrales : pour des
courbes elliptiques \emph{en famille}, au-dessus d'une base quelconque, où les
coefficients de l'équation sont des fonctions sur un espace et où même $2$ et
$3$ peuvent ne pas être inversibles. Il faut alors comprendre d'où vient
l'équation, ce qui la rend unique, et ce qu'on perd en la choisissant.

Deux idées y répondent. La première : l'équation n'est pas un choix arbitraire,
on la \emph{lit} sur la courbe. Les fonctions qui ont au plus un pôle d'ordre
$n$ en l'origine forment un espace de dimension $n$ ; une base adaptée de
l'espace de dimension $3$ plonge la courbe dans le plan et fournit l'équation ;
et toute la liberté dans ce choix est un groupe de matrices triangulaires. La
« variété des modules » est donc l'espace des équations divisé par ce groupe.
Une image : pour classer les triangles à déplacement près, on regarde les
triplets de sommets et l'on divise par le groupe des déplacements --- mais un
triangle isocèle a une symétrie, et le quotient n'est plus un espace ordinaire.
La page qui le dit (« pro-représentable », relativement aux « morphismes simples
surjectifs ») est l'endroit où l'on voit poindre ce qu'on appelle aujourd'hui
un champ.

La seconde idée est celle de Legendre. Une courbe elliptique est un revêtement
double de la droite, ramifié en quatre points ; on envoie trois d'entre eux en
$0$, $1$, $\infty$, et le quatrième devient un nombre $t$. Permuter les trois
premiers change $t$ en $1-t$, $1/t$, \ldots\ : six valeurs, et la fonction de
$t$ qui ne voit pas ces permutations est l'invariant $j$, dont les pages 27 à
39 calculent la formule explicite. Au-dessus d'une base, les trois points
peuvent s'échanger quand on se déplace : c'est un revêtement galoisien de
groupe $\mathfrak{S}_3$, et les pages 41 à 53 le traitent comme tel. On y lit,
en marge d'une même page, « fonctoriel » à côté d'une description valable sur
toute base, et « non fonctoriel » à côté d'une description par coordonnées,
valable seulement sur un anneau local --- le critère par lequel ces pages jugent
leurs propres énoncés.

Les noms à chercher ensuite : équations de Weierstrass sur une base, champ
des modules $\mathcal{M}_{1,1}$, fibré de Hodge et formes modulaires de poids
$4$, $6$, $12$, famille de Legendre et structures de niveau $2$, groupe
anharmonique, invariant $J$ de Klein, fonctions thêta et leur transformation
modulaire.

\keywords{elliptic curve over a base, Riemann-Roch theorem, Weierstrass equation, Lie algebra of an elliptic curve, Hodge bundle, moduli stack of elliptic curves, quotient stack, discriminant, Brauer-Severi scheme, cross-ratio, Legendre family, anharmonic group, j-invariant, level-2 structure, Galois cover, torsor, Kummer theory, Weierstrass elliptic functions, Weierstrass zeta function, Weierstrass sigma function, Legendre relation, Hermite decomposition, Jacobi elliptic functions, Jacobi imaginary transformation, Jacobi theta functions, heat equation, Klein J-invariant, modular group, elliptic points}
\end{resume}

\pagerange{1}{96}
\section*{Le fil du dossier, et ses conventions}

Le dossier a deux moitiés, séparées par une chemise (page 58), et la seconde se
subdivise en trois couches. Aucune page n'est numérotée par lui ; l'ordre est
celui des feuillets.

\begin{enumerate}
\item \emph{Pages 1 à 16} (chemise : « Courbes elliptiques » ; titre :
« Modules des courbes elliptiques ») --- les faisceaux $\mathcal{E}_n$ des
fonctions à pôles bornés le long de la section unité, le plongement dans le
plan, l'équation de Weierstrass sur une base, le schéma des équations et le
groupe triangulaire, les normalisations successives ; hors de la
caractéristique $2$, la courbe comme revêtement double d'une conique.
\item \emph{Pages 18 à 24} --- des feuilles de calcul : l'invariant $j$ d'une
cubique en fonction de ses coefficients et de ses racines, le discriminant, le
birapport. Les pages 19, 21 et 23 sont des versos dactylographiés sans rapport
(un tapuscrit numéroté « n° 276 », d'algèbre commutative), dont deux portent
une note au crayon.
\item \emph{Pages 27 à 39} --- la droite projective munie des trois points
$0, 1, \infty$ : son groupe d'automorphismes $\mathfrak{S}_3$, une involution
supplémentaire, les quotients successifs, et la formule explicite de $j$ en
fonction du paramètre de Legendre $t$.
\item \emph{Pages 41 à 53} --- la même chose « tordue » au-dessus d'une base :
une conique munie d'une trisection étale, sa classification par
$H^1(S, \mathfrak{S}_3)$, l'invariant $j$, la description par un
$\mathfrak{S}_3$-torseur ; puis « Autre façon de procéder ».
\item \emph{Pages 56 et 57} --- deux feuilles de brouillon : les poids des
coefficients de Weierstrass, un argument de descente.
\item \emph{Pages 58 à 96} (chemise : « Théorie transcendante (Généralités
classiques) ») --- la théorie analytique sur $\mathbf{C}$, en trois couches qui
se recouvrent et qu'on garde distinctes --- pages 59--60, d'une écriture posée
et d'aspect plus ancien : définition des fonctions elliptiques et théorèmes de
Liouville I à V ; pages 61--73, une première rédaction : $\wp$, $\zeta$,
$\sigma$, décomposition d'Hermite, formule de Jacobi, formules d'addition, puis
la fonction $\mathrm{sn}$ de Legendre et ses compagnes, reprises une seconde
fois sous les noms $\mathrm{sn}_1$, $\mathrm{cn}_1$, $\mathrm{dn}_1$ ; pages
74--96, une seconde rédaction, numérotée 1) à 5) : le corps des fonctions
elliptiques, $\wp$ et ses relations, $\zeta$, $\sigma$, les formules
d'addition, les fonctions $\sqrt{\wp - e_i}$, puis les fonctions thêta, la
fonction modulaire, et les fonctions de Legendre.
\end{enumerate}

Les deux moitiés se répondent sans se citer. L'équation $y^2 = 4x^3 - g_2 x -
g_3$ que la seconde moitié obtient pour $x = \wp$, $y = \wp'$ est celle que les
pages 8 à 16 extraient, sur une base quelconque, des faisceaux $\mathcal{E}_2$
et $\mathcal{E}_3$ : $\wp$ est la fonction à pôle double en l'origine, $\wp'$
la fonction à pôle triple. Les trois valeurs $e_1, e_2, e_3$ de $\wp$ aux
demi-périodes sont les trois points de ramification autres que l'origine du
revêtement double des pages 15 et 51 --- les points de $2$-torsion ---, et la
fonction $\mathrm{sn}^2 = 1/(\wp - e_1)$ de la page 66 est une coordonnée sur la
droite quotient qui envoie l'origine et deux d'entre eux en $0, \infty, 1$, et le
troisième en $1/k^2$ : à l'action de $\mathfrak{S}_3$ près, le module $k^2$ de
Legendre est le paramètre $t$ des pages 24 à 49. Enfin la fonction modulaire $J$ de la page 91 et l'invariant $j$ des
pages 39 à 47 sont liés par $j = 1 - J$.

\emph{Conventions, valables pour tout le dossier.} $S$ est un schéma
localement noethérien, $\pi\colon C \to S$ une courbe elliptique, $e$ sa section
unité et $D = e(S)$ le diviseur qu'elle définit.\footnote{Il note $f$ le
morphisme structural et $s$ la section unité ; on réserve $s$ à la section de
la conique des pages 43 à 49, $f$ au premier quotient de la page 35, et $D$ au
diviseur (il note aussi $D$ la tangente de la page 8 et la dérivation de la
page 13 ; on les appelle $T$ et $\partial$).} $\mathcal{L} =
\operatorname{Lie}(C/S) = e^*T_{C/S}$ est le faisceau de Lie (sa notation), et
$\omega = \mathcal{L}^\vee = e^*\Omega^1_{C/S}$ le fibré de Hodge.
$\mathcal{E}_n = \pi_*\mathcal{O}_C(nD)$, et $\operatorname{gr}_n\mathcal{E} =
\mathcal{E}_n/\mathcal{E}_{n-1}$.\footnote{Il écrit $\underline{L}$ pour
$\mathcal{O}_C(D)$ et $\mathcal{L}_n$ pour $\operatorname{gr}_n$ ; on évite la
seconde notation, qui se confondrait avec les puissances $\mathcal{L}^{\otimes n}$
--- auxquelles $\mathcal{L}_n$ est d'ailleurs isomorphe.} $x$ désigne une
fonction à pôle d'ordre $2$ et $y$ une fonction à pôle d'ordre $3$, selon
l'usage moderne.\footnote{Aux pages 8 à 16 il fait l'inverse : son $x$ est
d'ordre $3$, son $y$ d'ordre $2$, et son équation s'écrit $x^2 + ay^3 +
\cdots = 0$. Aux pages 24 et 56 il adopte déjà l'usage moderne. On l'adopte
partout, pour que $x$ et $y$ aient un seul sens dans le document.} Pour un
faisceau localement libre $\mathcal{F}$, $\mathbb{V}(\mathcal{F}) =
\operatorname{Spec}\operatorname{Sym}\mathcal{F}$ et $\mathbf{P}(\mathcal{F}) =
\operatorname{Proj}\operatorname{Sym}\mathcal{F}$, comme dans EGA --- c'est sa
convention.

L'invariant $j$ des pages 18 à 53 est le sien, et il n'est pas l'invariant
usuel : avec $j_{\mathrm{cl}} = 1728\,\dfrac{4a^3}{4a^3 + 27b^2}$ pour $y^2 =
x^3 + ax + b$, et $J = j_{\mathrm{cl}}/1728$ l'invariant de Klein de la page 91,
on a
\[
j = 1 - J = 1 - \frac{j_{\mathrm{cl}}}{1728} = \frac{27b^2}{4a^3+27b^2}.
\]
Ainsi $j = 0$ est la courbe à automorphismes d'ordre $4$ ($j_{\mathrm{cl}} =
1728$), $j = 1$ celle à automorphismes d'ordre $6$ ($j_{\mathrm{cl}} = 0$), et
$j = \infty$ la dégénérescence.\footnote{Cette identification est la nôtre ;
elle résulte de la formule (11) de la page 39 et de celles des pages 20 et 22,
et on l'a vérifiée par le calcul. La page ne compare jamais son $j$ à
l'invariant usuel.}

Dans la seconde moitié, $\mathfrak{P}$ est le réseau des périodes (sa lettre
gothique), engendré par $2\omega_1$ et $2\omega_2$ ; les demi-périodes sont
normalisées par $\omega_1 + \omega_2 + \omega_3 = 0$ (pages 62 et 90), $e_i =
\wp(\omega_i)$ et $\eta_i = \zeta(\omega_i)$. L'involution $t \mapsto
(t-2)/(2t-1)$ de la page 27 est notée $\kappa$, pour laisser $\sigma$ à la
fonction de Weierstrass.\footnote{Il la note $\sigma$.}

\emph{Ce que le dossier annonce et n'établit pas} : la démonstration de la
Proposition 3 (pages 6--7, bloc annulé) ; la forme explicite du discriminant
(page 10) ; la normalisation que demande la dernière ligne de la page 16 ; le
calcul du normalisateur de $\mathfrak{S}_3$ (page 29, « Prouver ») ; les données
supplémentaires qu'il faut quand $j$ n'est pas inversible (page 45) ; la fin de
la page 53 ; les quotients par $\kappa$ (page 39). On les retrouve à leur place.
\emph{Ce que la transcription ne donne pas} : la page 72, qu'annonce l'en-tête
de sa transcription (une troisième récapitulation de $\mathrm{sn}$, formules (1)
à (6)) mais que son texte ne contient pas ; la présente lecture ne la couvre
donc pas.

\pagerange{1}{16}
\section*{I. Modules des courbes elliptiques (pages 1 à 16)}

Le feuillet de garde porte « Courbes elliptiques », et la première page écrite
le titre « Modules des courbes elliptiques ».\footnote{Le titre de la chemise
est inscrit de sa main en haut à droite de la page 1, qui ne porte rien
d'autre ; la page 2 est blanche.}

\pagerange{3}{7}
\subsection*{Les faisceaux $\mathcal{E}_n$ (pages 3 à 7)}

La section unité $e$ définit un diviseur de Cartier relatif $D$, transversal
aux fibres, d'où le faisceau inversible $\mathcal{O}_C(D)$, dont les sections
sont les fonctions sur $C$ à pôle au plus simple le long de $D$. Sa restriction
à $D$ est le fibré normal, et celui-ci s'identifie à l'espace tangent le long
de la section unité :
\[
e^*\mathcal{O}_C(D) \simeq e^*\mathcal{N}_{D/C} \simeq e^*T_{C/S} =
\mathcal{L},
\qquad
e^*\mathcal{O}_C(nD) \simeq \mathcal{L}^{\otimes n}.
\]
Comme $C$ est un schéma en groupes, $T_{C/S} \simeq \pi^*\mathcal{L}$ : le
fibré tangent est trivialisé par translation.\footnote{La page commence
« $s^*(\underline{L}) \simeq$ » et le biffe pour poser d'abord $\underline{V}_{C/S}
= f^*(s^*\underline{V}_{C/S})$ ; une « Proposition 1 » commencée sous ces lignes est
barrée et reprise à la page 4. Le $S'$ de la seconde formule de la page est la
section unité vue comme sous-schéma de $C$.} La section $1$ de
$\mathcal{O}_C$ donne des injections $\mathcal{O}_C(nD) \hookrightarrow
\mathcal{O}_C((n+1)D)$ --- toutes contenues dans le faisceau des fonctions
méromorphes ---, d'où une suite croissante
\[
\cdots \subset \mathcal{E}_n \subset \mathcal{E}_{n+1} \subset \cdots,
\qquad \mathcal{E}_n = \pi_*\mathcal{O}_C(nD).
\]

\emph{Proposition 1.} Les faisceaux $\mathcal{O}_C(nD)$ sont
cohomologiquement plats sur $S$ en toute dimension, et
$\mathcal{E}_n$ est localement libre de rang $0$ si $n < 0$, de rang $1$ si
$n = 0$, de rang $n$ si $n \geqslant 1$. Le faisceau $\mathcal{O}_C(nD)$ est
ample si et seulement si $n \geqslant 1$, très ample si et seulement si $n
\geqslant 3$.

C'est Riemann--Roch en genre $1$, où $2g - 2 = 0$ : sur chaque fibre, un
faisceau inversible de degré $d \neq 0$ a $h^0 = d$, $h^1 = 0$ si $d > 0$, et
$h^0 = 0$, $h^1 = -d$ si $d < 0$ ; ces dimensions ne sautent pas, d'où la
platitude cohomologique, et pour $n = 0$ on a $\pi_*\mathcal{O}_C =
\mathcal{O}_S$, $R^1\pi_*\mathcal{O}_C \simeq \omega^\vee$, localement libres.
Le degré de $\mathcal{O}_C(nD)$ est $n$, d'où l'amplitude exactement pour $n
\geqslant 1$. Pour $n = 1, 2$, $\mathbf{P}(\mathcal{E}_n)$ est de dimension
relative $0$ ou $1$, et $C$ ne peut s'y plonger (il lui serait isomorphe) ;
pour $n \geqslant 3 = 2g + 1$, le critère usuel issu de Riemann--Roch donne la
très ample.\footnote{La démonstration de la page est en grande partie
illisible ; seuls la référence à Riemann--Roch, l'argument par la dimension de
$\mathbf{P}(\mathcal{E}_n)$ et l'inégalité « $3 > (2g-2)+2$ » se lisent. Une
note marginale au crayon parle d'une « prop.\ préliminaire ». L'argument
donné est le standard.}

\emph{Proposition 2.} $\mathcal{E}_0 = \mathcal{E}_1 = \mathcal{O}_S$, et pour
$i \leqslant j$ l'inclusion $\mathcal{E}_i \subset \mathcal{E}_j$ est
universellement injective --- localement facteur direct. Les quotients
$\operatorname{gr}_n\mathcal{E}$ sont nuls pour $n = 1$, égaux à
$\mathcal{O}_S$ pour $n = 0$, et
\[
\operatorname{gr}_n\mathcal{E} \simeq \mathcal{L}^{\otimes n}
\qquad (n \geqslant 2).
\]
En effet, la suite exacte $0 \to \mathcal{O}_C(nD) \to \mathcal{O}_C((n+1)D)
\to e_*\mathcal{L}^{\otimes(n+1)} \to 0$ donne
\[
0 \to \mathcal{E}_n \to \mathcal{E}_{n+1} \to \mathcal{L}^{\otimes(n+1)} \to
R^1\pi_*\mathcal{O}_C(nD),
\]
et le dernier terme est nul pour $n \geqslant 1$.\footnote{La page encadre
« $\mathcal{L}_{n+1} \simeq \mathcal{L}^{\otimes(n+1)}$, $n \geqslant 2$ » ;
l'argument vaut dès $n \geqslant 1$, c'est-à-dire pour $\operatorname{gr}_m$,
$m \geqslant 2$, ce qu'on énonce. Une note marginale demande d'« intervertir
avec la prop.\ 2 » : cette suite exacte, écrite à la page 7, est la bonne
démonstration de l'énoncé que la page 5 laissait en fragments.} Par suite la
multiplication induit des isomorphismes $\operatorname{gr}_n \otimes
\operatorname{gr}_m \simeq \operatorname{gr}_{n+m}$ pour $n, m \geqslant 0$,
$n, m \neq 1$, et l'algèbre graduée associée à la filtration est
\[
\operatorname{gr}\mathcal{E} = \mathcal{O}_S \oplus \bigoplus_{n \geqslant 2}
\mathcal{L}^{\otimes n}.
\]
Elle est engendrée par ses composantes de degrés $2$ et $3$ : pour $n
\geqslant 2$, $\operatorname{gr}_n$ est isomorphe à
$\operatorname{gr}_2^{\otimes n'}$ ou à $\operatorname{gr}_2^{\otimes n'}
\otimes \operatorname{gr}_3$ selon la parité de $n$ (Corollaire). Il en résulte
que l'algèbre $\mathcal{E}_\infty = \varinjlim \mathcal{E}_n =
\pi_*\mathcal{O}_C(*D)$ des fonctions à pôles le long de $D$ est engendrée par
$\mathcal{E}_3$, et de même l'anneau gradué $\bigoplus_k \mathcal{E}_{3k}$ de
$\mathcal{O}_C(3D)$.\footnote{La page dit « $\mathcal{E}_3$ engendre
$\Gamma_*(f, \underline{L})$ », où $\Gamma_* = \coprod_n \mathcal{E}_n$ est
l'anneau gradué de $\mathcal{O}_C(D)$ : pris à la lettre, c'est faux (les
degrés $1$ et $2$ ne sont pas engendrés par le degré $3$). Ce qui est vrai est
l'énoncé donné, filtré ou pour l'anneau de Veronese.}

\emph{L'action de l'algèbre de Lie} (pages 6 et 7, en partie annulées). Un
champ de vecteurs invariant $\partial$, section de $\mathcal{L}$, agit sur les
fonctions ; il élève l'ordre du pôle le long de $D$ d'une unité, et sur les
termes dominants, il multiplie par $-n$ une fonction à pôle d'ordre $n$
(car $\partial(z^{-n}) = -nz^{-n-1}$ pour un paramètre local $z$ adapté). D'où
des applications
\[
\mathcal{L} \otimes \operatorname{gr}_n\mathcal{E} \to
\operatorname{gr}_{n+1}\mathcal{E}, \qquad n \geqslant 2,
\]
qui, au travers de $\operatorname{gr}_n \simeq \mathcal{L}^{\otimes n}$, sont la
multiplication par $-n$.

\emph{Proposition 3.} En un point de $S$ de caractéristique résiduelle
d'exposant $p$, cette application est un isomorphisme si et seulement si $p$
est premier à $n$ ; si $p\mathcal{O}_S = 0$ et $p \mid n$, elle est nulle
(Corollaire).\footnote{L'énoncé et son corollaire sont lisibles ; la
démonstration (« vérification locale » sur une dérivation et l'ordre d'un
pôle) l'est par fragments, et l'argument donné ici est le nôtre. Deux
corollaires intermédiaires --- en caractéristique nulle $\mathcal{L}_n \simeq
\mathcal{L}^{\otimes(n-2)} \otimes \mathcal{L}_2$, en caractéristique $\neq 2$
$\mathcal{L}_2 \simeq \mathcal{L}^{\otimes 2}$ --- sont encadrés ensemble et
barrés : l'isomorphisme $\operatorname{gr}_n \simeq \mathcal{L}^{\otimes n}$
de la page 7, qui vaut en toute caractéristique, les rend inutiles.}

\pagerange{8}{10}
\subsection*{Le plongement cubique et l'équation de Weierstrass (pages 8 à 10)}

Par la Proposition 1, $\mathcal{O}_C(3D)$ définit une immersion fermée
\[
C \hookrightarrow \mathbf{P}(\mathcal{E}_3),
\]
dans un fibré en plans projectifs. L'image est un diviseur relatif de degré
$3$, défini par une forme cubique déterminée à une unité près : c'est le noyau,
localement libre de rang $1$, de $\operatorname{Sym}^3\mathcal{E}_3 \to
\mathcal{E}_9$ (rangs $10$ et $9$) --- un point de
$\mathbf{P}\bigl((\operatorname{Sym}^3\mathcal{E}_3)^\vee\bigr)$, que la page
écrit $\mathbf{P}(\operatorname{Symm}_3(\check{\mathcal{E}}_3))$.\footnote{La
description par le noyau est la nôtre ; la page dit seulement que $C$ est
« défini par une section sur $S$ de l'espace projectif des $3$-formes ».} Le
drapeau $\mathcal{E}_1 \subset \mathcal{E}_2 \subset \mathcal{E}_3$ a un sens
géométrique : $\mathbf{P}(\mathcal{E}_3/\mathcal{E}_2)$ est la section unité,
et $T = \mathbf{P}(\mathcal{E}_3/\mathcal{E}_1)$ est une droite relative, la
tangente en l'origine, qui est une \emph{tangente d'inflexion} --- elle coupe
$C$ en l'origine avec multiplicité $3$.\footnote{En marge, un croquis d'une
courbe et de sa tangente d'inflexion.}

Choisissons une base de $\mathcal{E}_3$ adaptée au drapeau et commençant par
$1$ : $1, x, y$, avec $x \in \mathcal{E}_2$, $y \in \mathcal{E}_3$. Deux telles
bases diffèrent par une matrice triangulaire
\[
\begin{pmatrix} 1 & \alpha & \gamma \\ 0 & \lambda & \beta \\ 0 & 0 & \mu
\end{pmatrix},
\qquad \lambda, \mu \in \mathcal{O}_S^*,
\]
qui dépend de cinq paramètres. Alors $1, x, y, x^2, xy, x^3$ est une base de
$\mathcal{E}_6$ adaptée à sa filtration, et $y^2 \in \mathcal{E}_6$ s'y écrit :
on obtient une équation
\[
\boxed{\;y^2 + a x^3 + b xy + c x^2 + d y + e x + f = 0\;}
\qquad (a \in \Gamma(S, \mathcal{O}_S^*)),
\]
où $a$ est inversible parce que $y^2$ et $x^3$ ont tous deux un terme dominant
non nul dans $\operatorname{gr}_6 \simeq \mathcal{L}^{\otimes 6}$.\footnote{Sur la
page, avec ses lettres, $x^2 + ay^3 + bxy + cy^2 + dx + ey + f = 0$ ; le
signe devant $ay^3$ est surchargé. Un premier essai, encadré et annulé, écrit
la forme générale à sept coefficients et un exposant illisible.} Homogénéisée
par une troisième coordonnée $z$, c'est la forme cubique
\[
\Phi(x, y, z) = y^2z + ax^3 + bxyz + cx^2z + dyz^2 + exz^2 + fz^3,
\]
et $\Phi$ \emph{est} l'équation de $C$ dans $\mathbf{P}(\mathcal{E}_3) \simeq
\mathbf{P}^2_S$ : si $\Phi'$ définit $C$, $\Phi$ est un multiple de $\Phi'$, et
le facteur est une unité puisque le coefficient de $y^2z$ vaut $1$ en tout
point de $S$.\footnote{La forme homogène de la page omet le terme en $yz^2$
(son $eyz^2$, notre $dyz^2$) ; on le rétablit. La page écrit
« $\simeq \mathbf{P}^3_S$ », exposant douteux à la lecture : c'est un plan.}

Réciproquement, des coefficients $a, \ldots, f \in \Gamma(S, \mathcal{O}_S)$
définissent un diviseur relatif positif de $\mathbf{P}^2_S$, qui passe par $(0
: 1 : 0)$ avec la tangente d'inflexion $z = 0$ ; il est lisse sur $S$ si et
seulement si un certain polynôme à coefficients entiers $\delta(a, \ldots, f)$,
le discriminant, est inversible en tout point.\footnote{La page écrit
« $\delta(a, b, c, d, e, f) \neq 0$ » ; deux lignes plus haut elle exige que
$\delta$ soit « une \emph{unité} partout », et c'est ce qu'on énonce. Elle
annonce « Ici le discriminant s'explicite : \ldots » et s'arrête : le polynôme
n'est pas écrit. Pour l'équation normalisée de la page 13 ($a = -1$), c'est le
discriminant $\Delta$ des formules de Tate, publiées par Deligne
(\emph{Courbes elliptiques : formulaire d'après J.~Tate}, 1975). Lisse
entraîne $a$ inversible : si $a$ s'annule en un point, la droite $z = 0$ est
contenue dans la fibre.} D'où :

\emph{Proposition} (page 10). Se donner une courbe elliptique sur $S$ munie
d'une « base spéciale » de $\mathcal{E}_3(C/S)$ --- adaptée au drapeau et
commençant par $1$ --- équivaut à se donner des sections $a, \ldots, f$ de
$\mathcal{O}_S$ avec $\delta(a, \ldots, f)$ inversible ; la correspondance est
compatible au changement de base.\footnote{Il avait d'abord écrit « base
adaptée », biffé pour « base spéciale ». Un couple (courbe, base spéciale) n'a
pas d'automorphisme non trivial, puisque la courbe est plongée par la base :
c'est pourquoi l'énoncé est une équivalence avec un \emph{ensemble} de
coefficients.}

\pagerange{10}{12}
\subsection*{Le schéma des équations et le groupe triangulaire (pages 10 à 12)}

Soit
\[
\mathfrak{M} = \operatorname{Spec}\bigl(\mathbf{Z}[A, B, C, D, E, F][\delta^{-1}]\bigr),
\]
ouvert affine de l'espace affine $\mathbf{E}^6_{\mathbf{Z}}$. Le groupe $G
\subset \mathrm{GL}(3)_{\mathbf{Z}}$ des matrices triangulaires ci-dessus opère
sur $\mathfrak{M}$ par substitution affine
\[
x' = \lambda x + \alpha, \qquad y' = \mu y + \beta x + \gamma
\]
dans la forme $y^2 + Ax^3 + Bxy + Cx^2 + Dy + Ex + F$.\footnote{Sur la page,
avec ses lettres, $Y' = \lambda Y + \alpha$, $X' = \mu X + \beta Y + \gamma$, et
la forme $X^2 + AY^3 + BXY^2 + \cdots$ ; l'exposant de $Y$ dans $BXY^2$ est une
lecture douteuse, et la page 9 a $bxy$, qui est seul de degré $5$.} La
Proposition de la page 10 dit que le foncteur « courbes elliptiques munies
d'une base spéciale » est \emph{représentable} par $\mathfrak{M}$. Deux points
de $\mathfrak{M}(S)$ définissent des courbes isomorphes si et seulement si ils
se déduisent l'un de l'autre, localement sur $S$, par un élément de $G$ ; et
toute courbe elliptique $C/S$ provient d'un point de $\mathfrak{M}(S')$ après
le changement de base $S' \to S$ par le $G$-torseur des bases spéciales de
$\mathcal{E}_3$. Ce torseur est lisse, surjectif, de type fini --- il n'est pas
fini.\footnote{La page a d'abord écrit « fid.\ plate », l'a biffé, et ajouté au
dessus « \ldots\ de type fini », précédé d'un mot lu « finie » sous réserve.
Le torseur sous $G$ est affine et lisse de dimension relative $5$.}

D'où la conclusion de la page 12 : le foncteur « classes d'isomorphisme de
courbes elliptiques sur $S$ » est « représentable » par le quotient de
$\mathfrak{M}$ par $G$, \emph{relativement à la catégorie des morphismes
simples surjectifs} --- c'est-à-dire après localisation pour la topologie
engendrée par les morphismes lisses surjectifs. En termes d'aujourd'hui : le
champ des modules des courbes elliptiques est le champ quotient\footnote{« Simple » est le mot qu'il emploie au début des années 1960 pour
« lisse ». Le mot de la page est lu « pro-représentable » sous réserve, et la
lettre gothique $\mathfrak{D}$ qui désigne la catégorie est écrite sur un mot
biffé ; la fin du paragraphe est presque entièrement illisible. La notion de
champ algébrique est postérieure (Deligne--Mumford, 1969 ; Artin, 1974) ; on
ne dit pas que la page la contient, seulement qu'elle décrit l'objet par une
présentation et une localisation, comme on le fait aujourd'hui. La présentation
usuelle, par les équations de Weierstrass $y^2 + a_1xy + a_3y = x^3 + a_2x^2 +
a_4x + a_6$ et le groupe des $(u, r, s, t)$, est celle de la page 13.}
\[
\mathcal{M}_{1,1} \simeq [\mathfrak{M}/G].
\]


\pagerange{13}{14}
\subsection*{Normalisation par une trivialisation de l'algèbre de Lie (pages 13 et 14)}

Une normalisation plus précise : on se donne une trivialisation $\partial$ de
$\mathcal{L}$ --- une dérivation invariante partout non nulle --- et l'on exige
que les images de $x$ et $y$ dans $\operatorname{gr}_2 \simeq
\mathcal{L}^{\otimes 2}$ et $\operatorname{gr}_3 \simeq \mathcal{L}^{\otimes 3}$
soient $\partial^{\otimes 2}$ et $\partial^{\otimes 3}$.\footnote{La page écrit
deux fois « $x$ » : « l'image de $x$ dans $\mathcal{L}_2$ soit $i_2(D)$, celle
de $x$ dans $\mathcal{L}_3$ soit $i_3(D)$ » ; la seconde est l'autre fonction
(notre $y$).} C'est-à-dire qu'on prend pour structure auxiliaire un faisceau
inversible $\mathcal{L}$ trivialisé et une base de $\mathcal{E}_3$ dont les
termes dominants sont $1, \partial^2, \partial^3$. Le groupe d'automorphismes
de cette structure est formé des couples
\[
\left(u,\ \begin{pmatrix} 1 & r & t \\ 0 & u^2 & s \\ 0 & 0 & u^3 \end{pmatrix}\right),
\qquad u \in \Gamma(S, \mathcal{O}_S^*),\ r, s, t \in \Gamma(S, \mathcal{O}_S),
\]
isomorphe au sous-groupe des matrices triangulaires de diagonale $(1, \lambda,
\mu)$ avec $\mu^2 = \lambda^3$ --- on retrouve $u = \mu/\lambda$.\footnote{Ses
lettres sont $(\lambda, \alpha, \beta, \gamma)$ ; on prend celles de Tate, $(u,
r, s, t)$ à l'ordre des coefficients près. Le cadre de la page porte
« $\mu^2 - \lambda^3 = 0$ », exposants surchargés.} Pour une telle base, $y^2 -
x^3 \in \mathcal{E}_5$, et l'équation prend la forme de Weierstrass
\[
y^2 + a_1xy + a_3y = x^3 + a_2x^2 + a_4x + a_6,
\]
c'est-à-dire $a = -1$ dans la forme précédente.\footnote{Avec ses lettres :
« $x^2 - y^3 + bxy + cy^2 + dx + ey + f = 0$ », « on peut prendre $a = -1$ ».}
Changer $\partial$ en $u\partial$ multiplie $a_i$ par $u^i$ : les coefficients
sont de poids $i$, et le discriminant $\Delta$, qui ne dépend pas de la base
une fois $\partial$ fixé, est une section canonique de $\omega^{\otimes 12}$ ---
une forme modulaire de poids $12$.\footnote{En marge de la page 9, entouré :
« $\Delta(a, b, c, d, e, f) \in \mathcal{L}^{\otimes 12}$ ». Avec la convention
suivie ici (un coefficient multiplié par $u^i$ quand $\partial$ l'est par $u$
est une section de $\mathcal{L}^{\otimes -i} = \omega^{\otimes i}$), c'est
$\omega^{\otimes 12} = \mathcal{L}^{\otimes -12}$ ; la page 56 a le signe que
l'on suit (« $b \in \Delta^{-4}$, $c \in \Delta^{-6}$ », où son $\Delta$ y est
le faisceau de Lie).}

La page 14 observe qu'on ne peut, localement, normaliser davantage sans
restreindre les courbes ; mais après une extension finie et plate on peut
exiger $\Delta = 1$, en choisissant la trivialisation $\partial$ parmi les
racines douzièmes de $\Delta^{-1}$, et pousser plus loin hors de la
caractéristique $2$, puis hors de la caractéristique $3$.\footnote{Le passage
est dense et lu par fragments ; « $12$-ième » est écrit sur un $11$, et une
parenthèse ajoute « ($24$-ième) », sans que la raison se lise. Que
$\omega^{\otimes 12}$ soit trivialisé par $\Delta$ est ce qui fait de
$\mathbf{Z}/12$ le groupe de Picard du champ des modules (Mumford,
\emph{Picard groups of moduli problems}, 1965) ; la page n'y renvoie pas, et
rien ne dit lequel des deux textes précède l'autre.}

\pagerange{15}{16}
\subsection*{Hors de la caractéristique 2 : le revêtement double (pages 15 et 16)}

L'inversion $-1$ de la courbe, qui est un schéma abélien sur $S$, est
d'ordre $2$ ; elle opère sur les $\mathcal{E}_i$ de façon compatible à la
multiplication, et induit sur $\mathcal{L}$ le changement de signe, donc $(-1)^n$
sur $\operatorname{gr}_n \simeq \mathcal{L}^{\otimes n}$. Supposons les
caractéristiques résiduelles de $S$ différentes de $2$. Alors $\mathcal{E}_3$
se décompose en parties paire et impaire : la partie paire est
$\mathcal{E}_2$, et la partie impaire est un supplémentaire canonique de
$\mathcal{E}_2$ dans $\mathcal{E}_3$, isomorphe à
$\operatorname{gr}_3$.\footnote{La page 16 est très raturée (deux lignes et
demie barrées, un passage encadré avec « À compléter », une note marginale
d'une dizaine de lignes lue par fragments) ; c'est la partie lisible, et elle
suffit. Il n'y a pas de supplémentaire canonique de $\mathcal{E}_1$ dans
$\mathcal{E}_2$ par la symétrie, puisque $\mathcal{E}_2$ est tout entier pair.}
Les bases « admissibles » sont celles où $y$ est impaire ; deux d'entre elles
diffèrent par une matrice
\[
\begin{pmatrix} 1 & r & 0 \\ 0 & u^2 & 0 \\ 0 & 0 & u^3 \end{pmatrix},
\]
et l'équation, comparée à sa transformée par $y \mapsto -y$, perd ses termes
impairs en $y$ :
\[
\boxed{\;y^2 = x^3 + ax^2 + bx + c\;}
\]
avec pour condition que le polynôme $P(x) = x^3 + ax^2 + bx + c$ ait des racines
distinctes en tout point de $S$.\footnote{La page fait le calcul sur
l'équation $x^2 = y^3 + axy + by^2 + cx + dy + e$ (ses lettres) et la symétrie
$x \mapsto -x$, dont l'introduction n'est pas sur la page. Elle finit par la
question « Quelle est la normalisation correspondante de l'équation de $C$ ? »
sans y répondre. Réponse : c'est l'équation encadrée, à $x \mapsto u^2 x + r$,
$y \mapsto u^3 y$ près ; si de plus $3$ est inversible, la translation $x
\mapsto x - a/3$ (calculée en marge de la page 20) tue $a$, et il reste $y^2 =
x^3 + bx + c$, à $(b, c) \mapsto (u^4b, u^6c)$ près, ce qui est la page 56.}

Géométriquement : le quotient $X = C/\{\pm 1\}$ est un schéma de
Brauer--Severi de dimension relative $1$ sur $S$, muni d'une section (l'image de
l'origine), donc une conique triviale ; privé de cette section, c'est un
torseur sous le groupe vectoriel $\mathbb{V}(\mathcal{L}^{\otimes 2})$ --- la
coordonnée $x$, définie à $x \mapsto x + r$ près et de terme dominant dans
$\operatorname{gr}_2 = \mathcal{L}^{\otimes 2}$, en est une fonction affine.
Et $C \to X$ est un revêtement double ramifié le long de la section à l'infini
et d'un diviseur étale de degré $3$ de $X$, contenu dans la partie affine : le
lieu $P(x) = 0$, image des points d'ordre $2$ non nuls.\footnote{La page écrit
« $C/\mathbf{Z}_2$ » pour le quotient, et « il correspond à
$\underline{L}^{\otimes 2}$ » pour le torseur. La dernière précision --- les
points d'ordre $2$ --- est la nôtre ; c'est elle qui fait le lien avec les
pages 24 à 53.}

\pagerange{18}{24}
\section*{II. Feuilles de calcul : l'invariant d'une cubique (pages 18 à 24)}

Quatre feuilles de calcul, sans rédaction, et trois versos
dactylographiés.\footnote{Les pages 19, 21 et 23 sont des feuillets d'un
tapuscrit numéroté « n° 276 » (anneaux $\mathfrak{m}$-adiques, modules filtrés
complets ; pages dactylographiées 20, 9 et 10), sans rapport avec les courbes
elliptiques. La page 19 ne porte rien de sa main. Les pages 21 et 23 portent
au crayon, dans la marge, une remarque chacune (« fait général \ldots\
métrisables ! » ; « Dire \ldots\ topologie \ldots\ coefficients »), d'une main
qui paraît être la sienne sans certitude. Il en va de même des pages 42, 44,
46, 48, 50, 52 et 54 plus loin, et des pages 25 à 40 paires.}

\pagerange{20}{20}
\subsection*{Le discriminant (page 20)}

Pour une cubique unitaire $P$ de racines
$x_1, x_2, x_3$,
\[
\prod_{i \neq j}(x_i - x_j) = \prod_i P'(x_i),
\]
et le discriminant usuel $\operatorname{disc}(P) = \prod_{i<j}(x_i - x_j)^2$ en
est l'opposé. Pour $P = x^3 + bx + c$ (la page suppose $a = 0$), le calcul de
$\prod(3x_i^2 + b)$ par les fonctions symétriques donne
\[
\prod_i P'(x_i) = 4b^3 + 27c^2, \qquad \operatorname{disc}(P) = -4b^3 - 27c^2,
\]
et son invariant est $j = -27c^2/\operatorname{disc}(P)$.\footnote{La page
encadre « $\Delta = b^3 - 27c^2$ » suivi de son propre point d'interrogation.
C'est le discriminant de Weierstrass $g_2^3 - 27g_3^2$ de la forme $4x^3 - g_2x
- g_3$ (pour $b = g_2$, $c = g_3$), pas celui de $x^3 + bx + c$ ; la formule $j
= -27c^2/\Delta$ est juste dans les deux normalisations, chacune avec son
$\Delta$. En marge, la translation $y \mapsto y - a/3$ qui tue le terme en
$y^2$.}

\pagerange{22}{22}
\subsection*{L'invariant en fonction des racines (page 22)}

Pour $P = x^3 + ax^2 + bx + c$, de
racines $x_i$ et de discriminant $\Delta$, on a $3x_i + a = 2x_i - x_j - x_k$,
et
\[
j = \frac{(27c + 2a^3 - 9ab)^2}{-27\,\Delta}
= -\frac{1}{27\Delta}\bigl[(2x_1 - x_2 - x_3)(2x_2 - x_3 - x_1)(2x_3 - x_1 -
x_2)\bigr]^2.
\]
En tête de page, l'identité qui relie ces expressions au paramètre $t$ des pages
24 et suivantes, $t = (x_3 - x_1)/(x_3 - x_2)$ :\footnote{La page écrit l'identité sans le signe $-$ ; on l'a vérifiée par le
calcul, elle vaut au signe près, ce qui est sans effet sur $j$, qui en est le
carré.}
\[
\frac{(t+1)(2t-1)(t-2)}{t(t-1)}
= -\,\frac{(2x_1 - x_2 - x_3)(2x_2 - x_3 - x_1)(2x_3 - x_1 - x_2)}{(x_1 -
x_2)(x_2 - x_3)(x_3 - x_1)} .
\]
Elle donne directement la formule (11) de la page 39, $j =
-\frac{1}{27}\bigl[(t+1)(2t-1)(t-2)/t(t-1)\bigr]^2$.

\pagerange{24}{24}
\subsection*{Forme normale et birapport (page 24)}

En tête, la forme normale
classique d'invariant donné,
\[
y^2 = 4x^3 - h(x+1), \qquad h = \frac{27\,j_{\mathrm{cl}}}{j_{\mathrm{cl}} - 2^6
3^3} = 27\,\frac{j-1}{j},
\]
la seconde expression étant celle de son invariant $j$ ; puis le birapport :
envoyant les racines $x_1, x_2, x_3$ et le point à l'infini sur $\infty, 0, 1,
t$, on trouve
\[
t = \frac{x_3 - x_1}{x_3 - x_2}, \qquad t - 1 = \frac{x_2 - x_1}{x_3 - x_2},
\qquad t + 1 = \frac{2x_3 - x_1 - x_2}{x_3 - x_2}, \ \ldots
\]
--- c'est le paramètre de Legendre, la quatrième valeur quand trois des quatre
points de ramification du revêtement double de la page 15 sont envoyés en
$\infty, 0, 1$.\footnote{La page porte encore des fractions en $a, b, c$
griffonnées au bord, la dernière coupée, et un croquis des quatre points sur
une droite.}

\pagerange{18}{18}
\subsection*{Quand $j$ ou $j - 1$ n'est pas inversible (page 18)}

La feuille
note, sur un anneau local $A$, ce qui se passe quand l'invariant réduit vaut $1$
ou $0$. Pour $y^2 = 4x^3 - bx - c$ ($b, c$ de poids $2$ et $3$ en
$\mathcal{L}$, selon le petit tableau de la page), on a
\[
\boxed{\;\frac{b^3}{c^2} = 27\,\frac{j-1}{j}\;}
\]
et si $j \equiv 1$, $b$ est nilpotent et $c$ inversible, la courbe n'étant
déterminée qu'à $c \in A^*/A^{*3}$ près ; si $j \equiv 0$, c'est $c$ qui est
nilpotent et $b$ inversible, à $b \in A^*/A^{*2}$ près. En marge : « Ext.\
Kummer de gr $\mathbf{Z}/3$ », « Ext.\ Kummer de gr $\mathbf{Z}/2$ ». Ce sont les
deux réductions du groupe structural de la page 47, et l'ambiguïté qui reste ---
un élément de $H^1(\operatorname{Spec}A, \mathbf{Z}/3)$ ou de
$H^1(\operatorname{Spec}A, \mathbf{Z}/2)$ --- est une extension de
Kummer.\footnote{Ce rapprochement est le nôtre, et la page ne le fait pas ; il
suppose qu'on oublie la torsion quadratique de la courbe, comme le fait la
donnée $(X, X', s)$ des pages 41 à 49. La feuille porte encore des racines
sixièmes $\sqrt[6]{t/(t-1)}$, une ligne « $(1 - j^{-1})^2 \in A^6A^*$ » et sa
compagne, lues en partie seulement, et sous le reste des
$\underline{\mathrm{Hom}}$ sans rapport apparent.} La page encadre aussi la
cubique $x^3 + 27\frac{j-1}{j}(x+1) = 0$ ; elle n'a pas l'invariant $j$, et la
forme normale juste est celle de la page 24, $4x^3 - 27\frac{j-1}{j}(x+1)$.\footnote{Pour
$x^3 + h(x+1)$, l'invariant vaut $27/(4h + 27)$, qui n'est $j$ que pour $h =
27(1-j)/4j$. Le calcul de la page repose sur le « $\Delta = b^3 - 27c^2$ » de la
page 20, juste pour $4x^3 - bx - c$, appliqué à une cubique unitaire.}

\pagerange{27}{39}
\section*{III. Le groupe anharmonique et la fonction $j$ (pages 27 à 39)}

\pagerange{27}{29}
\subsection*{Le groupe des automorphismes de $(\mathbf{P}^1, \{0, 1, \infty\})$ (pages 27 et 29)}

Sur $\mathbf{P}^1_S$, muni des trois sections $0, 1, \infty$, le groupe des
automorphismes qui respectent l'ensemble $\{0, 1, \infty\}$ s'identifie au groupe
$\mathfrak{S}_3$ de ses permutations --- un automorphisme de la droite
projective est déterminé par l'image de trois points distincts. Ses six
éléments sont
\[
1,\qquad \varphi(t) = 1 - \frac1t,\qquad \psi(t) = \frac{1}{1-t},\qquad
\lambda(t) = \frac1t,\qquad \mu(t) = 1 - t,\qquad \nu(t) = \frac{t}{t-1}:
\]
$\varphi$ et $\psi$ sont les deux permutations circulaires, $\lambda$, $\mu$,
$\nu$ les transpositions qui fixent respectivement $1$, $\infty$ et $0$. On note
$gh$ la substitution obtenue en appliquant d'abord $g$, puis $h$ ; avec cette
convention, la table de multiplication de la page est exacte\footnote{La convention n'est pas dite ; c'est celle qui rend la table exacte,
et on l'a vérifiée cas par cas. Pour la composition usuelle $g \circ h$,
chaque produit de deux éléments distincts serait remplacé par l'autre élément de
la même classe ($\lambda\circ\mu = \psi$, par exemple). En marge, une
proposition de renommer les transpositions, suivie de « ?? ».} :
\[
\varphi^2 = \psi,\quad \varphi\psi = 1,\quad \lambda^2 = \mu^2 = \nu^2 = 1,
\quad \lambda\mu = \mu\nu = \nu\lambda = \varphi,\quad \mu\lambda = \nu\mu =
\lambda\nu = \psi,
\]
\[
\varphi\lambda = \mu\varphi = \nu,\quad \varphi\mu = \nu\varphi = \lambda,\quad
\varphi\nu = \lambda\varphi = \mu,
\]
\[
\psi\nu = \mu\psi = \lambda,\quad
\psi\mu = \lambda\psi = \nu,\quad \psi\lambda = \nu\psi = \mu .
\]
Le sous-groupe
$\{1, \varphi, \psi\}$ est cyclique d'ordre $3$, distingué, de quotient
$\mathbf{Z}/2$ ; $\mathfrak{S}_3$ en est le produit semi-direct par l'un
quelconque des trois sous-groupes d'ordre $2$, engendrés par $\lambda$, $\mu$,
$\nu$. La conjugaison ${}^g h = g h g^{-1}$ (même convention) donne
\[
{}^\varphi\lambda = \mu,\ {}^\varphi\mu = \nu,\ {}^\varphi\nu = \lambda,\qquad
{}^\psi\nu = \mu,\ {}^\psi\mu = \lambda,\ {}^\psi\lambda = \nu,\qquad
{}^\tau\varphi = \psi\ \ (\tau = \lambda, \mu, \nu).
\]

De plus, l'involution
\[
\kappa(t) = \frac{t-2}{2t-1}, \qquad \kappa^2 = 1, \qquad \kappa(0) = 2,\
\kappa(1) = -1,\ \kappa(\infty) = \tfrac12,
\]
commute à $\mathfrak{S}_3$, si bien que le normalisateur de $\mathfrak{S}_3$
dans $\operatorname{Aut}(\mathbf{P}^1)$ contient $\mathfrak{S}_3 \times
\mathbf{Z}/2$ --- un groupe diédral d'ordre $12$ ; en caractéristique
résiduelle $\neq 2, 3$, sur un corps algébriquement clos, c'est tout le
normalisateur.\footnote{La page écrit « Normalisateur : $\mathfrak{S}_6$ (c'est
d'ailleurs le seul, \ldots » ; l'indice est peu net, et $\mathfrak{S}_6$ ne se
plonge pas dans $\mathrm{PGL}_2$. La note marginale de la page 29, « Prouver :
le normalisateur de $\mathfrak{S}_3$ dans $\operatorname{Aut}\mathbf{P}^1_S$ est
$\mathfrak{S}_3 \times (\mathbf{Z}/2)$ », dit ce qu'on énonce ; la preuve n'est
pas donnée. La commutation de $\kappa$ à $\mathfrak{S}_3$ est vérifiée ; la
dernière assertion vient de la classification des sous-groupes finis de
$\mathrm{PGL}_2$, que la page n'invoque pas. Un « 5 » appuyé en bout de ligne
n'a pas de rôle clair.}

\emph{Points fixes.} Les transpositions fixent chacune deux points :
$\lambda$ fixe $\{1, -1\}$, $\mu$ fixe $\{\infty, \frac12\}$, $\nu$ fixe $\{0,
2\}$ ; chaque paire est formée d'un point de $\{0, 1, \infty\}$ et d'un point
de $\kappa\{0, 1, \infty\} = \{2, -1, \frac12\}$. Les permutations circulaires
fixent les racines de
\[
t^2 - t + 1 = 0,
\]
c'est-à-dire de $\Phi_6$, les racines primitives sixièmes de l'unité $x' =
\frac12(1 + \sqrt{-3})$ et $x'' = \frac12(1 - \sqrt{-3})$ ; ce sont aussi les
points fixes de $\kappa$. Notons $\overline{S}$ ce sous-schéma de degré $2$ ;
il est étale sur $S$ exactement là où $3$ est inversible, et le quotient $\mathbf{Z}/2 = \mathfrak{S}_3/(\mathbf{Z}/3)$ y
échange les deux points.\footnote{La page écrit « $\mathbf{Z}_2$ » pour
$\mathbf{Z}/2$. Elle appelle « la différente » le lieu de ramification ; la
place de $\overline{S}$, écrit au-dessus de « l'un », n'est pas sûre dans la
phrase.}

\pagerange{29}{33}
\subsection*{Les quotients (pages 29 à 33)}

Le schéma $G$ des automorphismes de $\mathbf{P}^1$ qui laissent $\overline{S}$
invariant est une extension de $\mathbf{Z}/2$ (les automorphismes de
$\overline{S}$) par un $\mathbf{G}_m$ \emph{tordu} par le revêtement
$\overline{S} \to S$, et $\mathfrak{S}_3$ s'y envoie :\footnote{La page dessine les deux lignes dans l'ordre inverse, une seule
flèche verticale $\mathbf{Z}/3 \to \mathbf{G}_m$ et un signe $\|$ entre les
deux $\mathbf{Z}/2$ ; la flèche $\mathfrak{S}_3 \to G$ n'est pas tracée, on
l'ajoute parce que $\mathfrak{S}_3$ laisse $\overline{S}$ invariant. Une note
marginale renvoie aux points de caractéristique $2$ ou $3$ de $S$.}

\begin{tikzcd}
1 \arrow[r] & \mathbf{Z}/3 \arrow[r] \arrow[d, hook] & \mathfrak{S}_3 \arrow[r] \arrow[d, hook] & \mathbf{Z}/2 \arrow[r] \arrow[d, no head, "\Vert" description] & 1 \\
1 \arrow[r] & \mathbf{G}_m^{\mathrm{tordu}} \arrow[r] & G \arrow[r] & \mathbf{Z}/2 \arrow[r] & 1
\end{tikzcd}


On suppose désormais que $S$ n'a pas de point de caractéristique $2$ ou $3$.
Alors $\mathfrak{S}_3 \times \mathbf{Z}/2$ opère modérément sur
$\mathbf{P}^1_S$ (son ordre est inversible), le quotient est lisse sur $S$ de
genre $0$ --- c'est une forme de $\mathbf{P}^1_S$ ---, et il est
\emph{trivial} parce qu'il a des sections.\footnote{La page invoque Lüroth, qui
dit que les fibres sont des droites ; c'est l'existence des sections qui donne
l'isomorphisme global avec $\mathbf{P}^1_S$.} Les trois orbites
$\{0, 1, \infty\}$, $\{2, -1, \frac12\}$ et $\overline{S} = \{x', x''\}$ sont
disjointes, chacune stable sous $\mathfrak{S}_3$ qui y opère transitivement
fibre à fibre ; elles donnent dans $\mathbf{P}^1_S/\mathfrak{S}_3$ trois
sections distinctes, qu'on utilise pour identifier ce quotient à
$\mathbf{P}^1_S$ muni de $\infty, 0, 1$, dans cet ordre.

De même, $\mathbf{P}^1_S/(\mathbf{Z}/3)$ a deux sections, images de $\{0, 1,
\infty\}$ et de $\{2, -1, \frac12\}$, qui sont ses deux sections de
ramification, et qu'on envoie en $\infty$ et $0$ ; l'image de $\overline{S}$ est
une bisection étale. Sur $S' = S[\sqrt{-3}]$ elle se scinde en deux sections,
images de $x'$ et $x''$ ; en envoyant la première en $1$, la seconde va en
$-1$, et l'on obtient un isomorphisme $(\mathbf{P}^1_S/(\mathbf{Z}/3))
\times_S S' \simeq \mathbf{P}^1_{S'}$. Sur $S$ lui-même,
$\mathbf{P}^1_S/(\mathbf{Z}/3)$ est donc le fibré projectif
$\mathbf{P}(\mathcal{O}_S \oplus \mathcal{N})$ d'un faisceau inversible
$\mathcal{N}$ --- privé de $0$ et $\infty$, un fibré principal de groupe
$\mathbf{G}_{m,S}$ --- muni d'une bisection étale, ce qui revient à munir
$\mathcal{N}$ d'une forme quadratique non dégénérée $\mathcal{N}^{\otimes 2}
\simeq \mathcal{O}_S$ ; et la classe de cette donnée est celle du revêtement
$S' \to S$.\footnote{La page note ce faisceau $\Lambda_S$ ; on réserve
$\Lambda$ à la page 51. La dernière phrase de la page 33 se lit par
fragments.}

\pagerange{35}{39}
\subsection*{Formules explicites (pages 35 à 39)}

Notons $t$ la coordonnée de la première droite, $w$ celle de
$\mathbf{P}^1/(\mathbf{Z}/3)$ et $r$ celle de $\mathbf{P}^1/\mathfrak{S}_3$ :\footnote{Il note $s$ la coordonnée intermédiaire. Le schéma de la page
--- trois droites et les fibres marquées « ramifié » et « non ramifié » ---
est reproduit dans la liste qui suit ; ses deux fibres du haut sont écrites
sur des chiffres fortement surchargés.}
\[
\mathbf{P}^1_t \xrightarrow{\ f\ (\text{groupe } \mathbf{Z}/3)\ } \mathbf{P}^1_w
\xrightarrow{\ g\ (\text{groupe } \mathbf{Z}/2)\ } \mathbf{P}^1_r .
\]
Les fibres remarquables sont :
\begin{itemize}
\item au-dessus de $w = \infty$ et $w = 0$, les orbites $\{0, 1, \infty\}$ et
$\{2, -1, \frac12\}$, non ramifiées pour $f$ ; au-dessus de $w = 1$ et $w =
-1$, les points $x'$ et $x''$, ramifiés (indice $3$) ;
\item au-dessus de $r = \infty$ et $r = 0$, les points $w = \infty$ et $w =
0$, ramifiés pour $g$ ; au-dessus de $r = 1$, la paire $w = \pm 1$, non
ramifiée.
\end{itemize}

\emph{Calcul de $g$.} Sur $S = \operatorname{Spec}\mathbf{Z}[\frac12,
\frac13]$, la restriction de $g$ à $\mathbf{P}^1 \smallsetminus \{0, \infty\} =
\mathbf{G}_m$ est un endomorphisme de schémas $\mathbf{G}_m \to \mathbf{G}_m$
envoyant $1$ sur $1$, donc de la forme $w \mapsto cw^n$ ; $g(\infty) = \infty$
donne $n > 0$, le degré $2$ donne $n = 2$, et $g(1) = 1$ donne $c = 1$ :\footnote{Qu'un morphisme $\mathbf{G}_m \to \mathbf{G}_m$ soit $w \mapsto cw^n$
demande que les unités de $\mathcal{O}(S)[w, w^{-1}]$ soient les $cw^n$, ce
qui vaut sur une base réduite et connexe comme $\operatorname{Spec}\mathbf{Z}[\frac16]$.
La page dit « endomorphisme de $\mathbf{G}_m$ dans lui-même » ; une ligne sur
$\lambda$ de la forme $2^\alpha3^\beta$ est biffée. Elle ajoute que la formule
vaut sur tout $\mathbf{P}^1$, « et justifier cette notation ».}
\[
r = g(w) = w^2 . \tag{9}
\]


\emph{Calcul de $f$} (page 37). Sur $S'$, soit $u(t) = (t - x')/(t -
x'')$, qui envoie $x', x''$ sur $0, \infty$ ; dans cette coordonnée, $\mathbf{Z}/3$
opère par multiplication par les racines cubiques de l'unité, et le quotient est
$h(u) = -u^3$. Composant avec l'involution $v(w) = (1-w)/(1+w)$, qui
envoie $0, \infty, -1$ sur $1, -1, \infty$, on a $f = v^{-1} \circ h \circ u$ :

\begin{tikzcd}
\mathbf{P}^1_t \arrow[r, "u"] \arrow[d, "f"'] & \mathbf{P}^1 \arrow[d, "h"] \\
\mathbf{P}^1_w \arrow[r, "v"'] & \mathbf{P}^1
\end{tikzcd}

\[
f(t) = \frac{1 + u(t)^3}{1 - u(t)^3}
= \frac{(t - x'')^3 + (t - x')^3}{(t - x'')^3 - (t - x')^3}.
\]
Avec $x' + x'' = 1$, $x'x'' = 1$, $x'^3 = x''^3 = -1$, d'où $x'^2 + x''^2 = -1$,
$x'^2 + x'x'' + x''^2 = 0$ et $x' - x'' = \sqrt{-3}$, le numérateur vaut $2t^3
- 3t^2 - 3t + 2 = (t+1)(2t-1)(t-2)$ et le dénominateur $3\sqrt{-3}\,t(t-1)$ :\footnote{Le premier paragraphe de la page 37 est presque entièrement
illisible ; le calcul, lui, est complet, et on l'a refait. La page note
« $x''^3 - x'^3 = 0$ » pour la différence des cubes, ce qui est exact ; le
calcul du dénominateur passe par $x'^2 + x'x'' + x''^2 = 0$, que la page
n'écrit pas (une ligne en ce sens est biffée).}
\[
f(t) = \frac{1}{3\sqrt{-3}}\,\frac{(t+1)(2t-1)(t-2)}{t(t-1)} . \tag{10}
\]
On vérifie $f(0) = f(1) =
f(\infty) = \infty$, $f(-1) = f(2) = f(\frac12) = 0$, $f(x') = 1$.
Par suite
\[
\boxed{\;j = (g \circ f)(t) = -\frac{1}{27}\left(\frac{(t+1)(2t-1)(t-2)}{t(t-1)}\right)^2\;} \tag{11}
\]
est une fonction rationnelle de $t$, à coefficients dans $\mathbf{Z}[\frac16]$ :
la racine $\sqrt{-3}$ disparaît au carré. Elle définit un revêtement galoisien
de groupe $\mathfrak{S}_3$ de la droite des $j$, non ramifié au-dessus de
$\mathbf{P}^1 \smallsetminus \{0, 1, \infty\}$, c'est-à-dire au-dessus de
$\mathbf{G}_m \smallsetminus \{1\}$.\footnote{La page écrit « non ramifié
au-dessus de $(\mathbf{G}_a)_S - (0) - (1) = (\mathbf{G}_m)_S - 1$ » ; l'indice
$a$ est peu net. Le passage qui suit (11), « il doit définir \ldots\ sur $J$ !
\ldots\ $S$ lui-même », est en grande partie illisible.} En termes de
l'invariant usuel, $j_{\mathrm{cl}}(t) = 256\,(t^2 - t + 1)^3/t^2(t-1)^2$, et (11) dit
exactement $j = 1 - j_{\mathrm{cl}}/1728$ : c'est l'identité classique
$j_{\mathrm{cl}} - 1728 = 64\,(t+1)^2(2t-1)^2(t-2)^2/t^2(t-1)^2$.\footnote{Rapprochement
et vérification sont les nôtres.}

\emph{Les quotients de $\kappa$.} Comme $\kappa$ commute à $\mathfrak{S}_3$,
elle passe aux quotients ; elle échange les orbites au-dessus de $\infty$ et
de $0$, et fixe $x'$, d'où\footnote{La page écrit « $\sigma'(t) = \sigma''(t) = 1/t$ » ; l'identité
$j(\kappa t)\, j(t) = 1$ est vérifiée. Elle conclut qu'« il serait facile de
\ldots\ quotients \ldots », sans le faire.}
\[
\kappa'(w) = \frac1w, \qquad \kappa''(r) = \frac1r, \qquad\text{c'est-à-dire}\qquad
j(\kappa t) = 1/j(t). \tag{12}
\]


\pagerange{41}{53}
\section*{IV. Coniques munies d'une trisection étale (pages 41 à 53)}

\pagerange{41}{43}
\subsection*{Classification et invariant (pages 41 et 43)}

Soit $X$ un schéma de Brauer--Severi de dimension relative $1$ sur $S$ --- une
conique --- muni d'une \emph{trisection} étale $X'$ : un sous-schéma fermé de
$X$, fini et étale de degré $3$ sur $S$.\footnote{La page écrit deux fois « sur
$X$ » ; il faut $S$. Le mot en interligne, lu « fini », est
douteux.} Localement pour la topologie étale, $(X, X')$ est isomorphe à
$(\mathbf{P}^1_S, \{0, 1, \infty\})$, dont le groupe d'automorphismes est
$\mathfrak{S}_3$ (pages 27 à 29) ; ces couples sont donc classés par
$H^1(S, \mathfrak{S}_3)$, c'est-à-dire par les revêtements étales galoisiens
de groupe $\mathfrak{S}_3$ de $S$ --- ou, ce qui revient au même, par les
revêtements étales de degré $3$, dont $X'$ lui-même est un exemple.\footnote{La
dernière équivalence est la nôtre ; la page parle de « rev.\ non ramifiés
\ldots\ de groupe $\mathfrak{S}_3$ ».} Tout ce qui a été construit sur
$\mathbf{P}^1$ de façon invariante par $\mathfrak{S}_3$ se transporte à $X$
par descente : l'involution $\kappa$, qui envoie $X'$ sur une seconde
trisection $X''$, dite \emph{associée} à $X'$ ; et la bisection $X_1$, image
de $\overline{S}$, étale hors des caractéristiques $2$ et $3$.\footnote{La page
41 se lit par fragments ; « associé » (souligné) et « involutive » (en
interligne) sont lus sous réserve.}

Supposons $6$ inversible sur $S$. Le quotient $X/\mathfrak{S}_3$ --- le
groupe étant tordu comme $X$ --- est le quotient non tordu
$\mathbf{P}^1_S/\mathfrak{S}_3 \simeq \mathbf{P}^1_S$ de la page 31, et l'on
obtient un morphisme canonique de degré $6$\footnote{La page dit « isomorphisme unique $X \to \mathbf{P}^1_S$ caractérisé
\ldots\ par les sections $0, \infty, 1$ » : c'est un morphisme de degré $6$,
et c'est le quotient qui est isomorphe à $\mathbf{P}^1_S$. L'attribution des
trois fibres est en partie illisible et, telle qu'on la lit (« celle de
$\infty$ est $X''$ »), contredit la suite ; on suit ce qu'imposent la formule
(11), la phrase suivante (l'image d'une section disjointe de $X'$ est dans
$\mathbf{P}^1_S \smallsetminus \infty$) et la page 47 ($j = 0$ pour $X''$, $j =
1$ pour $X_1$).}
\[
q\colon X \to \mathbf{P}^1_S,
\qquad q^{-1}(\infty) = X',\quad q^{-1}(0) = X'',\quad q^{-1}(1) = X_1 .
\]


Donnons-nous de plus une section $s$ de $X$ disjointe de $X'$. Son image par $q$
est une section de $\mathbf{P}^1_S \smallsetminus \infty$, c'est-à-dire de
$\mathcal{O}_S$ : c'est l'\emph{invariant} $j$ du triple $(X, X', s)$. Si $j$
et $j - 1$ sont inversibles, l'image réciproque par $q$ de la section $j\colon
S \to \mathbf{P}^1_S$ est un revêtement étale galoisien de $S$ de groupe
$\mathfrak{S}_3$ --- $q$ est étale galoisien au-dessus de $\mathbf{P}^1
\smallsetminus \{0, 1, \infty\}$ --- et \emph{c'est celui qui définit $X$}.
Dans ce cas l'invariant suffit à reconstituer $(X, X', s)$.

Le lien avec les courbes elliptiques est celui de la page 15 : pour $C$ sur $S$
avec $6$ inversible, $X = C/\{\pm 1\}$, $X'$ l'image des points d'ordre $2$
non nuls et $s$ l'image de l'origine forment un tel triple, et dans une
coordonnée qui envoie $X'$ sur $0, 1, \infty$, $s$ est le paramètre $t$ de
Legendre (page 24). L'invariant $j$ du triple est alors $j(t)$, formule (11).
Le triple ne détermine la courbe qu'à torsion quadratique près : le revêtement
double ramifié le long de $s + X'$ demande le choix d'une racine carrée d'un
faisceau inversible.\footnote{Ce paragraphe est le nôtre ; le dossier ne
revient pas des triples aux courbes.}

\pagerange{45}{47}
\subsection*{Quand $j$ vaut $0$ ou $1$ (pages 45 à 47)}

Si $j$ n'est pas inversible, ou $j - 1$, l'invariant ne suffit plus ; quand
$j$ n'est pas identiquement $0$ ou $1$, on reconstitue $(X, X')$ sur l'ouvert
où $j(j-1)$ est inversible, puis $s$ partout, mais en général il faut des données
supplémentaires : la page demande d'exprimer qu'une certaine extension de
degré $3$ du corps des fonctions de $X$ est non ramifiée, et s'arrête.\footnote{La
page 45 est très lacunaire ; le passage « il faut des $H$ supplémentaires »
et la question « quelle \ldots\ convenablement ?? » sont lus sous réserve. La
page 46 est un verso du même tapuscrit « n° 276 » (sa page 15), où il a seulement
jeté de biais, à l'encre, les orbites « $(0, 1, \infty)$ », « $(x', x'')$ » et
« $(2, \frac12(2-1))$ ».} Dans les deux cas extrêmes, la situation se décrit
par une réduction du groupe structural :
\begin{itemize}
\item si $j = 0$, $s$ est dans $X''$, dont chaque point a pour stabilisateur
un sous-groupe d'ordre $2$ ; le groupe structural se réduit de $\mathfrak{S}_3$
à $\mathbf{Z}/2$ et la donnée est un élément de $H^1(S, \mathbf{Z}/2)$ ;
\item si $j = 1$, $s$ est dans la bisection $X_1$, de stabilisateur
$\mathbf{Z}/3$ ; le groupe se réduit à $\mathbf{Z}/3$ et la donnée est un
élément de $H^1(S, \mathbf{Z}/3)$.
\end{itemize}
Ce sont les extensions de Kummer de la page 18.

Si seulement $j$ est inversible, faisant opérer $\mathfrak{S}_3$ sur le
quotient intermédiaire $X/(\mathbf{Z}/3)$, les images de $X'$, $X''$ et $s$ y
donnent trois sections distinctes, qu'on envoie sur $\infty, 0, 1$ ; et $q$ se
factorise
\[
X \longrightarrow X/(\mathbf{Z}/3) \simeq \mathbf{P}^1_S
\xrightarrow{\ w \,\mapsto\, j w^2\ } \mathbf{P}^1_S,
\qquad (\infty, 0, 1) \mapsto (\infty, 0, j),
\]
la bisection naturelle s'identifiant aux points $w = \pm 1/\sqrt{j}$.\footnote{La
page écrit « $t \mapsto jt^2$ », la lettre $j$ lue sous réserve, et « la
bisection naturelle \ldots\ s'identifie donc \ldots\ : $\sqrt{j}$ » ; les
racines de $jw^2 = 1$ sont $\pm 1/\sqrt{j}$. La formule (9) est le cas $j = 1$
de cette normalisation.}

\pagerange{49}{49}
\subsection*{La description par un torseur (page 49)}

Ceci donne une description très explicite des courbes de genre $0$ sur $S$
munies d'une section et d'une trisection étale disjointes. Se donner $(X, X',
s)$ revient à se donner un $\mathfrak{S}_3$-torseur $T \to S$ et une fonction
\[
t \in \Gamma(T, \mathcal{O}_T), \qquad t \text{ et } t - 1 \text{ inversibles},
\qquad g\cdot t = g(t) \quad (g \in \mathfrak{S}_3),
\]
c'est-à-dire $\lambda\cdot t = 1/t$ et $\mu\cdot t = 1 - t$ : le torseur est
celui des isomorphismes $(X, X') \simeq (\mathbf{P}^1, \{0, 1, \infty\})$, et $t$
la coordonnée de $s$. Autrement dit, la catégorie de ces triples est celle des
$S$-points du champ quotient $[(\mathbf{P}^1 \smallsetminus \{0, 1,
\infty\})/\mathfrak{S}_3]$, et $X = T \times^{\mathfrak{S}_3}
\mathbf{P}^1_S$ :\footnote{Il note $S'$ le torseur (on garde $S'$ pour $S[\sqrt{-3}]$), et $\sigma
\in \mathfrak{S}_3$, $\varphi_\sigma(t)$ pour notre $g$, $g(t)$. Le champ
quotient est le nôtre. En marge, le long de ce paragraphe : « fonctoriel ».}

\begin{tikzcd}
X & \mathbf{P}^1_T \arrow[l] \arrow[d] \\
S & T \arrow[l, "\mathfrak{S}_3"]
\end{tikzcd}


\emph{Exemple.} Si $S = \operatorname{Spec}\mathcal{O}$, avec $\mathcal{O}$
local complet de caractéristique résiduelle $\neq 2, 3$ à corps résiduel
algébriquement clos, $\pi_1(S) = 0$ et tout torseur est trivial : la donnée
équivaut à celle d'un $t \in \mathcal{O}^*$ avec $t - 1 \in \mathcal{O}^*$, deux
tels $t$, $t'$ donnant des objets isomorphes si et seulement si $t' = g(t)$ pour
un $g \in \mathfrak{S}_3$. En marge, souligné deux fois : « non fonctoriel ».
Le contraste est le sien : la description par un torseur vaut sur toute base et
commute au changement de base ; celle par une coordonnée n'existe que là où les
torseurs sont triviaux, et l'identification dépend d'un choix.\footnote{La page
écrit « $t(t-1) \in \mathcal{O}^*$ » ; c'est la même condition. Un ajout
entouré, « Dém.\ \ldots », est illisible.}

\pagerange{51}{53}
\subsection*{Autre façon de procéder (pages 51 et 53)}

Prenons la section privilégiée $s$ comme section à l'infini. Alors $\Lambda = X
\smallsetminus s$ est un torseur sous le groupe vectoriel d'un faisceau
inversible --- pour $X = C/\{\pm 1\}$, celui de la page 15, $\mathbb{V}(\mathcal{L}^{\otimes 2})$.
La trisection $X'$ de $\Lambda$ a un \emph{barycentre} : dans un torseur sous un
groupe vectoriel, la moyenne de trois points est bien définie dès que $3$ est
inversible. Cette section de $\Lambda$ le trivialise, d'où un isomorphisme
\emph{global} de $\Lambda$ avec le groupe vectoriel lui-même, qui identifie $X'$
à une trisection de « somme » nulle. D'où l'énoncé de la page 53 :

\emph{Se donner $(X, X', s)$ équivaut à se donner un faisceau inversible
$\mathcal{M}$ sur $S$ et une trisection étale de $\mathbb{V}(\mathcal{M})$, de
somme nulle} ---
c'est-à-dire, dans une coordonnée locale, une cubique unitaire $x^3 + bx + c$ à
discriminant inversible, définie à $(b, c) \mapsto (u^4 b, u^6 c)$
près.\footnote{La dernière traduction est la nôtre ; elle rejoint la page 16
et la page 56. La page écrit « faisceau inversible $\underline{L}$ \ldots\
trisection $X'$ de $\mathcal{L} = V(\underline{L})$ », et « non ramifiée »,
entouré sous « inversible $\underline{L}$ », est relié par un trait à $X'$ : c'est
la trisection qui est non ramifiée. La phrase suivante, « Mais on peut aussi
éliminer les \ldots, symétriques \ldots\ trisection », et une note marginale
d'une dizaine de lignes, en partie biffée, ne se lisent pas.}

\pagerange{56}{57}
\section*{V. Deux feuilles de brouillon (pages 56 et 57)}

\emph{Les poids} (page 56). Pour une trivialisation $\xi$ de l'algèbre de Lie,
de dérivation invariante $D_\xi$, la fonction $u_\xi$ à pôle double vérifie
une équation de Weierstrass
\[
(D_\xi u_\xi)^2 = 4u_\xi^3 + b_\xi u_\xi + c_\xi .
\]
Remplacer $\xi$ par $\lambda\xi$ remplace $D_\xi$ par $\lambda D_\xi$ ; en posant
$u_{\lambda\xi} = \rho u_\xi$ et en identifiant, on trouve $\rho = \lambda^2$
et
\[
u_{\lambda\xi} = \lambda^2 u_\xi, \qquad b_{\lambda\xi} = \lambda^4 b_\xi,
\qquad c_{\lambda\xi} = \lambda^6 c_\xi :
\]
$b$ et $c$ sont des sections de $\omega^{\otimes 4}$ et $\omega^{\otimes 6}$
--- des formes modulaires de poids $4$ et $6$ --- et $u$ a son terme dominant
dans $\operatorname{gr}_2 \simeq \mathcal{L}^{\otimes 2}$.\footnote{La page note
$\Delta$ le faisceau de Lie (« $u_\xi \in L_2 \simeq \Delta^2$ »,
« $b \in \Delta^{-4}$, $c \in \Delta^{-6}$ »), après avoir écrit à tort
$\Delta^0$ pour $b_\xi$ et $c_\xi$ --- ce sont des fonctions une fois $\xi$
fixé, d'où sans doute ce $\Delta^0$. Deux relations plus petites, lues
« $\Delta^? \simeq \mathcal{L}^{18}$, $\Delta^2 \simeq \mathcal{L}^3$ », restent
obscures. La même feuille récrit la filtration $\mathcal{O} = E_1 \subset E_2
\subset \cdots$ avec ses quotients en puissances de $\Delta$, et note $y^2 = 4x^3
+ bx + c$, $y/x = D$.}

\emph{Un argument de descente} (page 57). Soit $A \to B$ un homomorphisme
d'anneaux, $A$ local d'idéal maximal $\mathfrak{m}$, et $x_i \in B$ tels que
les $x_i \otimes 1$ forment une base de $B \otimes_A B$ sur $B$ (agissant par le
second facteur). Par Nakayama, les $x_i$ engendrent $B$ sur $A$ ; et si
$\sum\lambda_i x_i = 0$ avec $\lambda_i \in A$, alors $\sum (x_i \otimes
1)(1 \otimes \lambda_i) = 0$, donc $1 \otimes \lambda_i = 0$, et par
$B \otimes_A B \to B$, $\lambda_i = 0$ dans $B$. Si $A \to B$ est injectif ---
par exemple fidèlement plat ---, les $x_i$ forment donc une base de $B$ sur
$A$.\footnote{La feuille est un brouillon sans phrase ; l'énoncé ainsi
reconstitué est le nôtre, et l'hypothèse d'injectivité, nécessaire pour
conclure de « $\lambda_i = 0$ dans $B$ » à « $\lambda_i = 0$ », n'est pas sur la
page. En tête, un diagramme $X \times_Y X \to X \to Y$ et une relation $R
\subset X \times_S X$ : le contexte est celui des quotients des pages 31 et 43.}

\pagerange{58}{60}
\section*{VI. Théorie transcendante : les théorèmes de Liouville (pages 58 à 60)}

Le feuillet de garde de la page 58 porte « Théorie transcendante (Généralités
classiques) ». Ce qui suit, jusqu'à la page 96, est la théorie des fonctions
elliptiques sur $\mathbf{C}$, en trois couches (sections VI à VIII) ; la
première tient en deux pages.\footnote{L'inscription est au crayon, de sa
main, en haut à droite de la page 58, qui ne porte rien d'autre. Les pages 59 et
60 sont sur papier brun, d'une encre brune et d'une écriture posée, plus
ancienne d'aspect ; les pages 61 à 73 sont d'une main régulière, les pages 74 à
96 d'une main plus rapide.}

Une \emph{fonction elliptique} est une fonction méromorphe sur $\mathbf{C}$
admettant deux périodes $\omega_1$, $\omega_2$ linéairement indépendantes sur
$\mathbf{Z}$. Si elle n'est pas constante, $\omega_1/\omega_2$ n'est pas
réel.\footnote{La page dit « deux périodes distinctes » ; il faut
l'indépendance sur $\mathbf{Z}$ ($\sin z$ a les périodes $2\pi$ et $4\pi$), et
c'est alors un théorème que le rapport n'est pas réel --- s'il était réel
irrationnel, les périodes seraient denses. Deux lignes ajoutées en interligne,
lues par fragments, demandent de « remplacer “$\omega_1$ et $\omega_2$” ».} Les
parallélogrammes de sommets $a + n_1\omega_1 + n_2\omega_2$ pavent le plan :
un point intérieur à l'un n'appartient qu'à lui, un point d'un côté (hors
sommet) à deux, un sommet à quatre ; deux points intérieurs d'un même
parallélogramme ne sont pas homologues. Quand aucun zéro ni pôle n'est sur
le bord, la liste des zéros et des pôles avec multiplicité ne dépend pas du
parallélogramme ; en général, on compte une seule fois chaque classe de points
homologues.\footnote{La page compte un point des côtés pour moitié et un
sommet pour un quart, ce qui revient au même. Les indices $n_1$, $n_2$ varient
d'une formule à l'autre sur la page.}

\begin{itemize}
\item[I.] Une fonction elliptique sans pôle est constante.
\item[II.] La somme des résidus d'une fonction elliptique dans un
parallélogramme des périodes est nulle ; une fonction elliptique non constante y
a donc au moins deux pôles, comptés avec multiplicité.
\item[III.] Le nombre des zéros et celui des pôles dans un parallélogramme sont
égaux, et égaux au nombre de fois que la fonction y prend une valeur donnée :
c'est l'\emph{ordre} de la fonction.
\item[--] Deux fonctions elliptiques de mêmes périodes qui ont mêmes zéros et
mêmes pôles avec multiplicité sont égales à un facteur constant près ; si elles
ont mêmes pôles avec mêmes parties principales, à une constante additive près.
\item[IV.] La somme des zéros est égale à la somme des pôles, à une période près.
\item[V.] Deux fonctions elliptiques de mêmes périodes sont liées par une
relation algébrique.\footnote{Le théorème I n'est pas numéroté sur la page (une croix dans la
marge) ; le chiffre romain du quatrième énoncé est sous une tache d'encre, et un
trait le relie au théorème II. Entre deux, la remarque que la somme, le
produit, le quotient et la dérivée de fonctions elliptiques de mêmes périodes
sont de mêmes périodes, et que l'ensemble des valeurs prises dans un
parallélogramme est celui des valeurs prises dans le plan.}
\end{itemize}

\pagerange{61}{73}
\section*{VII. Première rédaction : $\wp$, $\zeta$, $\sigma$, $\mathrm{sn}$ (pages 61 à 73)}

\pagerange{61}{64}
\subsection*{Les fonctions $\wp$, $\zeta$, $\sigma$ (pages 61 à 64)}

\emph{La fonction $\wp$} (page 61). C'est la fonction elliptique d'ordre $2$,
de périodes $2\omega_1$, $2\omega_2$, ayant en l'origine un pôle double et
telle que $\wp(z) - 1/z^2$ s'annule en $0$ ; ces conditions la déterminent. Elle
est paire --- $\wp(z) - \wp(-z)$ est elliptique, sans pôle, nulle en $0$ ---,
d'où
\[
\wp(z) = \frac{1}{z^2} + c_1z^2 + c_2z^4 + \cdots, \qquad
\wp'^2 = 4\wp^3 - g_2\wp - g_3, \quad g_2 = 20c_1,\ g_3 = 28c_2 ;
\]
elle est la fonction réciproque de $z = \int_\infty^{Z} dZ/\sqrt{4Z^3 - g_2Z -
g_3}$. Les valeurs $e_i = \wp(\omega_i)$ aux demi-périodes sont les racines de
$4Z^3 - g_2Z - g_3$, et $\wp'(\omega_i) = 0$. On a $\wp'' = 6\wp^2 -
\frac12 g_2$, $\wp''' = 12\wp\wp'$ ; $\wp^{(2n)}$ est un polynôme en $\wp$ de
degré $n+1$, et $\wp^{(2n+1)}$ le produit de $\wp'$ par un polynôme de degré $n$.

\emph{La fonction $\zeta$} (pages 61 et 62) :
\[
\zeta(z) = \frac1z - \int_0^z \Bigl(\wp(u) - \frac{1}{u^2}\Bigr)du,
\qquad \zeta' = -\wp ,
\]
primitive de $-\wp$, impaire, à pôles simples de résidu $1$ aux
périodes.\footnote{La page dit « une primitive de $\wp z$ » ; la formule qui
suit donne, correctement, $\zeta' = -\wp$.} Elle n'est pas périodique :
$\zeta(z + 2\omega_i) - \zeta(z) = 2\eta_i$ avec $\eta_i = \zeta(\omega_i)$, et
une intégration de $\zeta$ le long du bord d'un parallélogramme donne la
relation de Legendre\footnote{L'orientation n'est pas dite ici ; elle l'est à la page 80 (« si
$\omega_1$, $\omega_2$ sont en disposition directe »). Pour l'autre orientation
le second membre est $-i\pi/2$. La somme nulle des $\eta_i$ suppose $\omega_1 +
\omega_2 + \omega_3 = 0$, notre convention.}
\[
\eta_1\omega_2 - \eta_2\omega_1 = \frac{i\pi}{2}
\qquad \Bigl(\Im\frac{\omega_2}{\omega_1} > 0\Bigr),
\qquad \eta_1 + \eta_2 + \eta_3 = 0 .
\]
Elle permet d'intégrer $Z\,dZ /
\sqrt{4Z^3 - g_2Z - g_3}$ : avec $Z = \wp(z)$, $\int \wp(z)\,dz = -\zeta(z) +
\text{c}^{\text{te}}$.

\emph{Décomposition d'Hermite.} Si $f$ est elliptique, de pôles $a, b,
\ldots$ deux à deux non homologues, de parties principales $\sum_k
A_k/(z-a)^k$, $\sum_k B_k/(z-b)^k$, \ldots, alors
\[
f(z) = \text{c}^{\text{te}} + \sum_{\text{pôles } a}\Bigl( A_1\zeta(z-a) +
A_2\wp(z-a) - \frac{A_3}{2!}\wp'(z-a) + \cdots + (-1)^\alpha
\frac{A_\alpha}{(\alpha-1)!}\wp^{(\alpha-2)}(z-a)\Bigr),
\]
le second membre étant elliptique parce que $\sum A_1 = 0$ (théorème II).

\emph{La fonction $\sigma$} (pages 62 et 63). Pour intégrer toute fonction
elliptique, il reste à intégrer $\zeta$, ce qui amène $\log\sigma$. Comme
\[
\wp(z) = \frac{1}{z^2} + \sum_{w \in \mathfrak{P}\smallsetminus 0}
\Bigl[\frac{1}{(z-w)^2} - \frac{1}{w^2}\Bigr],
\qquad c_1 = 3\sum_{w \neq 0} \frac{1}{w^4},\quad c_2 = 5\sum_{w \neq 0}
\frac{1}{w^6},
\]
\[
\zeta(z) = \frac1z + \sum_{w \neq 0}\Bigl[\frac{1}{z-w} + \frac1w +
\frac{z}{w^2}\Bigr],
\qquad
\sigma(z) = z\prod_{w \neq 0}\Bigl(1 - \frac{z}{w}\Bigr)e^{\frac{z}{w} +
\frac{z^2}{2w^2}},
\]
le produit convergeant absolument et uniformément sur les compacts, $\sigma$
est une fonction entière, nulle exactement aux périodes, avec $\sigma'/\sigma =
\zeta$ ; c'est la fonction entière la plus simple qui s'annule en ces points. Elle
est \emph{impaire}, et\footnote{La page écrit « $\sigma z$ est une fonction paire » ; elle est
impaire, comme le dit d'ailleurs la page 81.}
\[
\sigma(z + 2\omega_i) = -e^{2\eta_i(z + \omega_i)}\sigma(z) .
\]


\emph{Formule de Jacobi} (page 64). Soit $f$ elliptique, de pôles $a, \ldots,
l$ et de zéros $a_1, \ldots, l_1$ (avec multiplicité, deux à deux non
homologues à une période près), et $a + \cdots + l = a_1 + \cdots + l_1 + \Omega$,
$\Omega = 2n_1\omega_1 + 2n_2\omega_2$ (théorème IV). Alors\footnote{La page écrit l'exponentielle $e^{2(n_1\eta_1 + n_2\eta_2)z}$ ; le
signe est $-$, comme le donne la quasi-périodicité de $\sigma$ appliquée à
$\sigma(z - l_1 - \Omega)$, et comme on l'a vérifié numériquement. Les deux
premières expressions de la page, dont la seconde déplace $\Omega$ au
dénominateur ($\sigma(z - l + \Omega)$), sont justes.}
\[
f(z) = \text{c}^{\text{te}} \times \frac{\sigma(z - a_1)\cdots\sigma(z - l_1 -
\Omega)}{\sigma(z-a)\cdots\sigma(z-l)}
= \text{c}^{\text{te}} \times e^{-2(n_1\eta_1 + n_2\eta_2)z}
\frac{\sigma(z-a_1)\cdots\sigma(z-l_1)}{\sigma(z-a)\cdots\sigma(z-l)} .
\]


\emph{Formules d'addition} :
\[
\wp u - \wp v = -\frac{\sigma(u+v)\sigma(u-v)}{\sigma^2u\,\sigma^2v},
\qquad
\zeta(u+v) = \zeta u + \zeta v + \frac12\,\frac{\wp'u - \wp'v}{\wp u - \wp v},
\]
\[
\wp(u+v) + \wp u + \wp v = \frac14\Bigl(\frac{\wp'u - \wp'v}{\wp u - \wp v}\Bigr)^2
= \bigl[\zeta(u+v) - \zeta u - \zeta v\bigr]^2 .
\]

\pagerange{65}{69}
\subsection*{La fonction $\mathrm{sn}$ de Legendre et ses compagnes (pages 65 à 69)}

L'intégrale
$\int dZ/\sqrt{(1-Z^2)(1-k^2Z^2)}$ est tentante par sa symétrie, et redonne les
fonctions trigonométriques et hyperboliques pour $k = 0$ et $k = 1$. On appelle
$\mathrm{sn}\,z$ la solution de
\[
\Bigl(\frac{dZ}{dz}\Bigr)^2 = (1 - Z^2)(1 - k^2Z^2), \qquad Z(0) = 0,\ Z'(0) = 1 .
\]
On l'exprime par une fonction $\wp$ de trois façons :
\begin{itemize}
\item[1°] par une homographie, $Z = \dfrac{12\wp(z - z_0) - (5 - k^2)}{12\wp(z -
z_0) - (5k^2 - 1)}$, pour le $\wp$ de racines $e = \frac16(1+k^2)$,
$-\frac{1}{12}(k^2 \pm 6k + 1)$, avec $\wp(z_0) = (5 - k^2)/12$ ; trop lourd,
dit la page, mais cela prouve l'existence ;\footnote{L'identité différentielle
est vérifiée : elle donne $\wp'(z_0)^2 = (1-k^2)^2/4$, d'où la valeur
$\wp'(z_0) = -(1-k^2)/2$ de la page, au signe près fixé par $Z'(0) = 1$.}
\item[2°] et 3°, pour le $\wp$ de racines
\[
e_1 = -\frac{1 + k^2}{3}, \qquad e_2 = \frac{2 - k^2}{3}, \qquad e_3 =
\frac{2k^2 - 1}{3},
\]
\[
\mathrm{sn}^2 z = \frac{1}{k^2}\bigl[\wp(z - \omega_1) - e_1\bigr], \qquad
\boxed{\;\mathrm{sn}^2 z = \frac{1}{\wp z - e_1}\;} .
\]
\end{itemize}
L'uniformité de $\mathrm{sn}$ vient de ce que les racines de $\wp z - e_1$ sont
doubles. Les deux expressions donnent l'identité $(\wp z - e_1)(\wp(z - \omega_1)
- e_1) = (e_1 - e_2)(e_1 - e_3) = k^2$, valable pour toute fonction $\wp$ (avec
$(e_1 - e_2)(e_1 - e_3)$ au second membre), et qui vaut aussi
\[
k^2 = \frac{e^{2\eta_1\omega_1}}{\sigma(\omega_1)^4} = 3e_1^2 - \frac{g_2}{4} .
\]
Ces normalisations imposent les demi-périodes : $\omega_1 = iK'$ (pôle de
$\mathrm{sn}$), $\omega_2 = K$, $\omega_3 = -K - iK'$, où $K$, $K'$ sont les
intégrales complètes de la page 93. On a alors
\[
\mathrm{sn}(z + \omega_1)\,\mathrm{sn}\,z = \frac1k,
\qquad
\mathrm{sn}\,z = \sigma(\omega_1)e^{\eta_1 z}\frac{\sigma z}{\sigma(z + \omega_1)}
= -\sigma(\omega_1)e^{-\eta_1 z}\frac{\sigma z}{\sigma(z - \omega_1)} ,
\]
\[
\mathrm{sn}(z + 2\omega_1) = \mathrm{sn}\,z, \quad \mathrm{sn}(z + 2\omega_2) =
\mathrm{sn}(z + 2\omega_3) = -\mathrm{sn}\,z, \quad \mathrm{sn}(-z) =
-\mathrm{sn}\,z :
\]
$\mathrm{sn}$ a les périodes $2\omega_1$ et $4\omega_2$, elle est d'ordre $2$,
ses zéros sont homologues de $0$ et $2\omega_2$ (dérivée $1$ et $-1$), ses pôles
homologues de $\omega_1$ et $\omega_1 + 2\omega_2$ (résidus $1/k$ et $-1/k$) ;
enfin $\mathrm{sn}\,\omega_2 = 1$, $\mathrm{sn}'\omega_2 = 0$, $\mathrm{sn}(\omega_1 +
\omega_2) = 1/k$, $\mathrm{sn}\,\omega_3 = -1/k$, $\mathrm{sn}'\omega_3 =
0$.\footnote{Toutes ces formules ont été vérifiées numériquement. Avec ces
étiquettes, $\Im(\omega_2/\omega_1) < 0$, et la relation de Legendre de la page
62 prend le signe $-i\pi/2$. Sur la page, les pôles sont écrits $\omega_1 +
2n_1\omega_1 + 4n_2\omega_2$ et les zéros $2n_1\omega_1 + 4n_2\omega_2$ --- ce qui
n'en donne que la moitié --- ; les zéros sont dits homologues de « $0$ ou
$2\omega_1$ » et les pôles de « $\omega_2$ et $\omega_2 + 2\omega_1$ » :
$\omega_1$ et $\omega_2$ y sont échangés, contre (3), (5), (6) et (7) de la même
page, qu'on suit. La normalisation du signe de $k$ est donnée sous la forme
« $k = -e^{\eta_1\omega_2}/\sigma'\omega_2$ », indices lus sous réserve ; ce qui
vaut est $k = \pm e^{\eta_1\omega_1}/\sigma(\omega_1)^2$, le carré de la formule
précédente. La valeur « $\mathrm{sn}\,\omega_2 = 1/?$ » a un dénominateur
illisible ; c'est $1$. Le raisonnement sur la parité, à la page 67, se lit mal ;
la conclusion, $\mathrm{sn}$ impaire, est juste.}

On pose (page 68)
\[
\mathrm{cn}^2 z = 1 - \mathrm{sn}^2 z = \frac{\wp z - e_2}{\wp z - e_1} =
-\frac{1}{k^2}\bigl[\wp(z - \omega_1) - e_3\bigr], \qquad \mathrm{cn}\,0 = 1,
\]
\[
\mathrm{dn}^2 z = 1 - k^2\mathrm{sn}^2 z = \frac{\wp z - e_3}{\wp z - e_1} =
-\bigl[\wp(z - \omega_1) - e_2\bigr], \qquad \mathrm{dn}\,0 = 1,
\]
\[
\mathrm{cn}\,z = \frac{\sigma\omega_1}{\sigma\omega_2}e^{(\eta_2 - \eta_1)z}
\frac{\sigma(z - \omega_2)}{\sigma(z - \omega_1)},
\qquad
\mathrm{dn}\,z = \frac{\sigma\omega_1}{\sigma\omega_3}e^{(\eta_3 - \eta_1)z}
\frac{\sigma(z - \omega_3)}{\sigma(z - \omega_1)},
\]
et les formes équivalentes en $z + \omega_i$. Les fonctions $\mathrm{cn}$ et
$\mathrm{dn}$ sont paires, d'ordre $2$, et
\[
\mathrm{cn}(z + 2\omega_1) = \mathrm{cn}(z + 2\omega_2) = -\mathrm{cn}\,z,\quad
\mathrm{cn}(z + 2\omega_3) = \mathrm{cn}\,z ;
\]
\[
\mathrm{dn}(z + 2\omega_2) = \mathrm{dn}\,z,\quad \mathrm{dn}(z + 2\omega_1) =
\mathrm{dn}(z + 2\omega_3) = -\mathrm{dn}\,z ,
\]
si bien que $\mathrm{cn}$ a les périodes $2\omega_3$, $4\omega_1$, $4\omega_2$, et
$\mathrm{dn}$ les périodes $2\omega_2$, $4\omega_1$, $4\omega_3$. Les pôles de
$\mathrm{cn}$ et de $\mathrm{dn}$ sont homologues de $\pm\omega_1$, de résidus
$\mp i/k$ pour $\mathrm{cn}$ et $\mp i$ pour $\mathrm{dn}$ ; les zéros de
$\mathrm{cn}$ sont homologues de $\pm\omega_2$, avec $\mathrm{cn}'\omega_2 = -k'$,
ceux de $\mathrm{dn}$ de $\pm\omega_3$, avec $\mathrm{dn}'\omega_3 = -ik'$, où
$k' = \sqrt{1 - k^2}$ ; enfin $\mathrm{cn}\,\omega_3 = -ik'/k$, $\mathrm{dn}\,\omega_2 =
k'$, $\mathrm{cn}'\omega_3 = \mathrm{dn}'\omega_2 = 0$.\footnote{Vérifié
numériquement. La page donne $\mathrm{cn}(z + 2\omega_2) = \mathrm{cn}\,z$ et
$\mathrm{cn}(z + 2\omega_3) = -\mathrm{cn}\,z$, d'où les périodes « $2\omega_2$,
$4\omega_1$, $4\omega_3$ » pour $\mathrm{cn}$ et « $2\omega_3$, $4\omega_1$,
$4\omega_2$ » pour $\mathrm{dn}$ : $\omega_2$ et $\omega_3$ y sont échangés, et la
transcription signale plusieurs indices mal lisibles dans ces lignes. La
quasi-périodicité de $\sigma$ appliquée à la formule de $\mathrm{cn}$ par
$\sigma$, qui est sur la page, donne ce qu'on écrit. Les pôles de $\mathrm{cn}$
et de $\mathrm{dn}$ sont dits homologues de $\omega_2$ (et les zéros de
$\mathrm{cn}$ aussi) ; il faut $\omega_1$. Les valeurs y sont données au signe
près ($\pm\sqrt{1-k^2}$, $\pm\frac{i}{k}\sqrt{1-k^2}$), avec la remarque que
l'ambiguïté se lève par les formules par $\sigma$ ; on donne les signes. La
relation « $\mathrm{cn}'^2\omega_2 = -\mathrm{sn}'^2\omega_2$ » ne peut être
juste telle quelle ($\mathrm{sn}'\omega_2 = 0$) et n'est pas reprise.}

Les formules de dérivation et les relations quadratiques :
\[
\mathrm{sn}' = \mathrm{cn}\,\mathrm{dn},\quad \mathrm{cn}' =
-\mathrm{sn}\,\mathrm{dn},\quad \mathrm{dn}' = -k^2\,\mathrm{sn}\,\mathrm{cn} ;
\]
\[
\mathrm{cn}^2 + \mathrm{sn}^2 = 1,\quad \mathrm{dn}^2 + k^2\mathrm{sn}^2 = 1,\quad
\mathrm{dn}^2 - k^2\mathrm{cn}^2 = 1 - k^2 ;
\]
\[
\mathrm{cn}'^2 = (1 - \mathrm{cn}^2)(k'^2 + k^2\mathrm{cn}^2),
\qquad
\mathrm{dn}'^2 = (1 - \mathrm{dn}^2)(\mathrm{dn}^2 - k'^2) .
\]
À un changement affine de la variable près, $\mathrm{cn}$ et $\mathrm{dn}$ sont
donc solutions de l'équation de $\mathrm{sn}$ pour les modules $\rho$ tels que
$\rho^2 = -k^2/(1-k^2)$ et $\rho^2 = 1/(1-k^2)$ respectivement.\footnote{La page
écrit $\rho = k/\sqrt{1-k^2}$ pour $\mathrm{cn}$ : il manque le facteur $i$, que
le signe $+$ dans $k'^2 + k^2\mathrm{cn}^2$ impose. Pour $\mathrm{dn}$, $\rho =
1/\sqrt{1-k^2}$ est juste. Le numéro « (12) » de la page se lit « (11) ».}

\pagerange{70}{73}
\subsection*{Seconde reprise : $\mathrm{sn}_1$, $\mathrm{cn}_1$, $\mathrm{dn}_1$ (pages 70 à 73)}

La page 70 reprend le tout pour une fonction $\wp$ quelconque, sans
normaliser $k$ :
\[
\mathrm{sn}_1^2 z = \frac{1}{\wp z - e_1}, \qquad \mathrm{sn}_1'(0) = 1,
\qquad
\mathrm{sn}_1(z)\,\mathrm{sn}_1(z + \omega_1) = \frac{1}{k_1},\quad
k_1^2 = (e_2 - e_1)(e_3 - e_1) = \frac{e^{2\eta_1\omega_1}}{\sigma(\omega_1)^4},
\]
avec la même formule par $\sigma$ ; $\mathrm{sn}_1$ a les périodes $2\omega_1$,
$4\omega_2$, $4\omega_3$, change de signe par $z \mapsto z + 2\omega_2$ ou $z
+ 2\omega_3$, est impaire et invariante par $z \mapsto 2\omega_2 - z$ et $z
\mapsto 2\omega_3 - z$ ; ses pôles sont homologues de $\omega_1$ (résidu
$1/k_1$) et de $\omega_1 + 2\omega_2$ (résidu $-1/k_1$), ses zéros de $0$
(dérivée $1$) et de $2\omega_2$ (dérivée $-1$) ; $\mathrm{sn}_1\omega_2 =
\pm 1/\sqrt{e_2 - e_1}$, $\mathrm{sn}_1\omega_3 = \pm 1/\sqrt{e_3 - e_1}$, et
$\mathrm{sn}_1'$ s'annule en ces deux points.\footnote{Le second facteur de
$k_1^2$ se lit « $(e_3 - e_1)$ » une seconde fois, un indice étant peu net ; la
normalisation est donnée par « $k_1 = -e^{\eta_1\omega_1}/\sigma'\omega_1$ », où
il faut lire $\sigma(\omega_1)^2$, comme plus haut. Dans la formule par
$\sigma$ de $\mathrm{sn}_1\omega_3$, coupée par le bas de la feuille, le
dénominateur se lit $\sigma(\omega_1 + \omega_2)$ ; c'est $\sigma(\omega_1 +
\omega_3)$. Ici les pôles sont bien homologues de $\omega_1$.} Puis
\[
\mathrm{cn}_1^2 = 1 - (e_2 - e_1)\,\mathrm{sn}_1^2 = \frac{\wp - e_2}{\wp - e_1},
\qquad
\mathrm{dn}_1^2 = 1 - (e_3 - e_1)\,\mathrm{sn}_1^2 = \frac{\wp - e_3}{\wp - e_1},
\]
\[
\mathrm{sn}_1' = \mathrm{cn}_1\mathrm{dn}_1,\quad \mathrm{cn}_1' = -(e_2 -
e_1)\,\mathrm{sn}_1\mathrm{dn}_1,\quad \mathrm{dn}_1' = -(e_3 -
e_1)\,\mathrm{sn}_1\mathrm{cn}_1,
\]
\[
\mathrm{sn}_1'^2 = \bigl(1 - (e_2 - e_1)\mathrm{sn}_1^2\bigr)\bigl(1 - (e_3 -
e_1)\mathrm{sn}_1^2\bigr).
\]
Pour les racines $e_i$ de la page 66 on a $e_2 - e_1 = 1$ et $e_3 - e_1 = k^2$, et
l'on retrouve $\mathrm{sn}$, $\mathrm{cn}$, $\mathrm{dn}$.\footnote{La page 71
rappelle sous un trait les valeurs « $\mathrm{sn}\,\omega_2 = 1$,
$\mathrm{sn}\,\omega_3 = -1/k$ », qui sont celles de la fonction normalisée.
Les formules de la page 73 sont numérotées (8) et (9), sans (7) ; la page 72,
qui portait selon l'en-tête de la transcription une récapitulation de
$\mathrm{sn}$ numérotée (1) à (6), n'est pas transcrite, et n'est donc pas lue
ici.}

\pagerange{74}{96}
\section*{VIII. Seconde rédaction : « Fonctions elliptiques » (pages 74 à 96)}

\pagerange{74}{78}
\subsection*{1) Généralités (pages 74 à 78)}

Une fonction elliptique est une fonction
méromorphe $f$ sur $\mathbf{C}$ admettant deux périodes linéairement
indépendantes sur $\mathbf{R}$ ; si $f$ n'est pas constante, son groupe des
périodes, sous-groupe fermé de $\mathbf{C} = \mathbf{R}^2$, est discret, donc un
réseau $\mathfrak{P}$. Fixons-le, avec une base $\{2\omega_1, 2\omega_2\}$
(définie à une substitution entière de déterminant $\pm 1$ près) : les fonctions
méromorphes vérifiant $f(z + 2\omega_1) = f(z + 2\omega_2) = f(z)$ forment un
corps $E$, celui des fonctions méromorphes sur la courbe analytique
$\mathbf{C}/\mathfrak{P}$, topologiquement un tore. La théorie de Riemann dit
que $E$ est de degré fini sur $\mathbf{C}(f)$ pour tout $f$ non constant ; on
va obtenir directement plus précis.\footnote{Le titre « Fonctions elliptiques
» est encadré, avec en sous-titre « Expression des fonctions elliptiques par
$\wp u$, $\wp' u$ ». Une phrase « Une fonction elliptique sans pôle est une
constante » et un début de « a) Existence \ldots » sont rayés.}

\begin{itemize}
\item[a$'$)] La somme des résidus est nulle, et le nombre des zéros est égal au
nombre des pôles, comptés avec multiplicité --- par le théorème des résidus sur
le bord d'un parallélogramme des périodes.\footnote{La page dit « immédiat par
Liouville », souligné. Liouville donne qu'une fonction elliptique sans pôle est
constante ; les deux énoncés de a$'$) viennent de l'intégration de $f$ et de
$f'/f$ sur le bord.}
\item[b)] La série
\[
\wp(z) = \frac{1}{z^2} + \sum_{(m,n) \neq (0,0)}\Bigl(\frac{1}{(z - 2m\omega_1 -
2n\omega_2)^2} - \frac{1}{(2m\omega_1 + 2n\omega_2)^2}\Bigr)
\]
définit une fonction méromorphe sur $\mathbf{C}$, à pôles doubles aux points
de $\mathfrak{P}$, et $\wp' = -2\sum_{m,n}(z - 2m\omega_1 - 2n\omega_2)^{-3}$ est
visiblement elliptique. De $\wp'(z) = \wp'(z + 2\omega_i)$ on tire $\wp(z + 2\omega_i)
- \wp(z) = \text{c}^{\text{te}}$, nulle comme on le voit en $z = -\omega_i$
puisque $\wp$ est paire. C'est l'unique fonction elliptique n'ayant de pôles
qu'aux points de $\mathfrak{P}$ et telle que $\wp(z) - 1/z^2$ soit holomorphe
et nulle en $0$ ; elle est paire, $\wp(z) = 1/z^2 + c_1z^2 + c_2z^4 +
\cdots$.\footnote{La page dit « holomorphe dans le plan » pour la série, et
« pôles \ldots\ $2\mathfrak{P}$ » semble-t-il : $\mathfrak{P}$, qui est déjà
l'ensemble des $2m\omega_1 + 2n\omega_2$.}
\item[c)] Toute fonction elliptique paire n'ayant de pôles qu'aux points de
$\mathfrak{P}$ est un polynôme en $\wp$, et d'une seule façon (car $\wp$ est
transcendante sur $\mathbf{C}$) ; si le pôle en $0$ est d'ordre $2k$, de terme
dominant $c_k z^{-2k}$, le polynôme est de degré $k$ et de coefficient dominant
$c_k$ : $f - c_k\wp^k$ est paire, à pôle d'ordre $\leqslant 2(k-1)$, et l'on
conclut par récurrence.\footnote{La page ajoute « Le polynôme est unitaire »,
qui est faux ($2\wp$) ; la démonstration qu'elle donne produit le coefficient
dominant $c_k$.}
\item[d)] Toute fonction elliptique paire est une fonction rationnelle de
$\wp$, et réciproquement. Si $a_1, \ldots, a_n$ sont ses pôles hors de
$\mathfrak{P}$, un par paire $\{a, -a\}$ modulo $\mathfrak{P}$, d'ordres $l_i$,
la fonction $f\prod_i(\wp - \wp(a_i))^{l_i}$ est paire et n'a de pôles que dans
$\mathfrak{P}$ ; on applique c).\footnote{Cette démonstration est celle de la
page 77. La page 76 est entièrement barrée de longs traits obliques : elle
tentait une autre voie (le dénominateur de $R_2$ dans $\wp'R_2(\wp)$, et un
lemme : « une fonction elliptique a au moins deux pôles non congruents mod
$\mathfrak{P}$, ou un pôle double »), et établissait que $\wp'$, impaire, a
pour zéros simples les trois demi-périodes --- ce que la page 78 reprend.}
\item[e)] $[E : \mathbf{C}(\wp)] = 2$ : $\mathbf{C}(\wp)$ est, par d), le corps
des invariants de l'automorphisme $f \mapsto \check f$, $\check f(z) = f(-z)$,
qui est d'ordre $2$ --- et $E \neq \mathbf{C}(\wp)$ puisque $\wp'$ est impaire.
\item[f)] Toute fonction elliptique s'écrit d'une seule façon
\[
f = R_1(\wp) + \wp' R_2(\wp), \qquad R_1, R_2 \text{ rationnelles} ;
\]
\item[g)] ses pôles sont dans $\mathfrak{P}$ si et seulement si $R_1$ et $R_2$
sont des polynômes ; une telle fonction, non constante, a en $0$ un pôle au
moins double. Une fonction impaire est de la forme $\wp' R(\wp)$, avec $R$
polynôme si et seulement si ses pôles sont dans $\mathfrak{P}$.\footnote{Une
preuve de g) par $f = \frac12(f + \check f) + \frac12(f - \check f)$ est
encadrée et barrée à la page 77 ; la page 78 donne la réduction au cas impair,
écrite vite.}
\end{itemize}

\pagerange{78}{81}
\subsection*{2) Les relations entre $\wp$ et $\wp'$ ; 3) la fonction $\zeta$ (pages 78 à 81)}

En dérivant la
série, les coefficients sont
\[
c_k = (2k+1)\sum_{w \in \mathfrak{P}\smallsetminus 0}\frac{1}{w^{2k+2}} .
\]
Les zéros de $\wp'$ dans $\mathbf{C}/\mathfrak{P}$ sont les trois demi-périodes
$\omega_1, \omega_2, \omega_3$ ; posant $e_i = \wp(\omega_i)$, ce sont des
racines doubles de $\wp z - e_i = 0$, et comme $\wp - e_i$ n'a que deux zéros
dans $\mathbf{C}/\mathfrak{P}$, \emph{les $e_i$ sont distincts}.\footnote{La page
écrit « dans $\mathfrak{P}$ » pour $\mathbf{C}/\mathfrak{P}$, et ici $\omega_3 =
\omega_1 + \omega_2$ : c'est le même point de $\mathbf{C}/\mathfrak{P}$ que
notre $-\omega_1 - \omega_2$, et $e_3$ ne change pas ; seul $\eta_3$ change de
signe.} Une fonction elliptique et sa dérivée étant liées par une relation
algébrique, toute fonction elliptique vérifie une équation différentielle
algébrique du premier ordre. Pour $\wp$ : $\wp'^2$ est paire, à pôle d'ordre $6$,
donc un polynôme de degré $3$ en $\wp$, et l'identification des développements
en $0$ donne
\[
\boxed{\;\wp'^2 = 4\wp^3 - g_2\wp - g_3\;}, \qquad g_2 = 20c_1,\ g_3 = 28c_2 ;
\]
ou encore : $\wp'^2/\prod(\wp - e_i)$ est elliptique sans pôle, donc constante,
égale à $4$, et\footnote{Les numéros des deux formules encadrées se lisent mal (un $4$ ou un
$7$, un $5$ ou un $8$) ; la page 80 continue par (6). Une ligne finale écrit le
polynôme « $4(\wp^3 - e_1)(\wp - e_2)(\wp - e_3)$ » ; c'est $\wp - e_1$.}
\[
\boxed{\;\wp'^2 = 4(\wp - e_1)(\wp - e_2)(\wp - e_3)\;} .
\]
En
dérivant, $\wp'' = 6\wp^2 - \frac12 g_2$ et $\wp''' = 12\wp\wp'$, et en poussant
l'identification,
\[
c_3 = \frac{c_1^2}{3}, \qquad c_4 = \frac{3}{11}c_1c_2 .
\]

\emph{3) La fonction $\zeta$} (pages 80 et 81). $\zeta(u) - \frac1u = -\int_0^u
(\wp z - \frac{1}{z^2})dz$, $\zeta' = -\wp$,
\[
\zeta(u) = \frac1u - \frac{c_1}{3}u^3 - \cdots - \frac{c_n}{2n+1}u^{2n+1} -
\cdots,
\qquad
\zeta(z) = \frac1z + \sum_{w \neq 0}\Bigl(\frac{1}{z - w} + \frac1w +
\frac{z}{w^2}\Bigr) ;
\]
$\zeta$ est impaire, à pôles simples de résidu $1$ aux points de $\mathfrak{P}$.
Comme $\zeta'$ est elliptique, $\zeta(z + 2\omega) - \zeta(z) = 2\eta$ est
constant pour toute demi-période $\omega$, et en $z = -\omega$ on trouve $\eta
= \zeta(\omega)$. Les $\eta_i$ et les $\omega_i$ sont liés par la relation de
Legendre, $\eta_1\omega_2 - \eta_2\omega_1 = i\pi/2$ quand $\omega_1, \omega_2$
sont en disposition directe ; la page 81 en indique la preuve par une
intégrale le long d'un parallélogramme des périodes.\footnote{Le premier
coefficient se lit $\frac{c_2}{3}$, indice surchargé ; c'est $\frac{c_1}{3}$, le
terme général de la même ligne le dit.}

\emph{Formule d'Hermite} (page 81). Comme à la page 62, pour $f$ de pôles $a_i$
et de parties principales $\sum_{m=1}^{k_i} A_{i,m}/(z - a_i)^m$,
\[
f(z) = \text{c}^{\text{te}} + \sum_i\Bigl(A_{i,1}\zeta(z - a_i) + A_{i,2}\wp(z -
a_i) + \cdots + \frac{(-1)^{k_i}}{(k_i - 1)!}A_{i,k_i}\wp^{(k_i - 2)}(z -
a_i)\Bigr), \tag{12}
\]
elliptique grâce à $\sum_i A_{i,1} = 0$.\footnote{L'ordre de dérivation de
$\wp$ se lit mal sur la page ; c'est $k_i - 2$.}

\pagerange{81}{84}
\subsection*{4) La fonction $\sigma$ ; 5) formules d'addition ; les fonctions $\wp - e_i$ (pages 81 à 84)}

\emph{4) La fonction $\sigma$} (pages 81 et 82).
\[
\sigma(z) = z\prod_{w \in \mathfrak{P}\smallsetminus 0}\Bigl(1 -
\frac{z}{w}\Bigr)e^{\frac{z}{w} + \frac{z^2}{2w^2}}, \qquad
\frac{\sigma'}{\sigma} = \zeta, \qquad \sigma(z) = z\,e^{\int_0^z(\zeta u -
\frac1u)du}, \tag{13--14}
\]
fonction entière impaire, à zéros simples exactement aux points de
$\mathfrak{P}$, et\footnote{La page écrit (16) $\sigma(u + 2\omega_i) = \sigma u\,e^{2\eta_i(u +
\omega_i) + 2i\pi}$ ; le facteur $e^{2i\pi} = 1$ perd le signe. Il faut $+i\pi$,
comme la page l'écrit elle-même en (26) et à la page 84.}
\[
\sigma(u) = u - \frac{c_1}{12}u^5 - \frac{c_2}{30}u^7 - \cdots, \qquad
\sigma(u + 2\omega_i) = -\sigma(u)\,e^{2\eta_i(u + \omega_i)} . \tag{15--16}
\]
\emph{Formule de
Jacobi} : si $(a_i)$ et $(b_i)$ sont des systèmes de représentants des zéros
et des pôles avec $\sum a_i = \sum b_i$ (ce que permet le théorème IV),
\[
f(u) = \text{c}^{\text{te}} \times \frac{\prod\sigma(u - a_i)}{\prod\sigma(u -
b_i)} , \tag{17}
\]
le second membre étant elliptique par (16).

\emph{5) Formules d'addition} (pages 82 et 83). La formule de Jacobi appliquée
à $u \mapsto \wp u - \wp v$ donne
\[
\wp u - \wp v = -\frac{\sigma(u+v)\sigma(u-v)}{\sigma^2u\,\sigma^2v} ; \tag{18}
\]
ses dérivées logarithmiques en $u$ et en $v$,
\[
\frac{\wp'u}{\wp u - \wp v} = \zeta(u+v) + \zeta(u-v) - 2\zeta u, \qquad
-\frac{\wp'v}{\wp u - \wp v} = \zeta(u+v) - \zeta(u-v) - 2\zeta v, \tag{19}
\]
\[
\zeta(u+v) - \zeta u - \zeta v = \frac12\,\frac{\wp'u - \wp'v}{\wp u - \wp v},
\qquad \zeta(2u) - 2\zeta u = \frac12\,\frac{\wp''u}{\wp'u}, \tag{20}
\]
\[
\wp(u+v) + \wp u + \wp v = \frac14\Bigl(\frac{\wp'u - \wp'v}{\wp u -
\wp v}\Bigr)^2, \qquad \wp(2u) + 2\wp u = \frac14\Bigl(\frac{\wp''u}{\wp'u}\Bigr)^2 ;
\tag{21}
\]
et la formule d'Hermite appliquée à $\wp(2u)$ :
\[
4\wp(2u) = \wp u + \wp(u + \omega_1) + \wp(u + \omega_2) + \wp(u + \omega_3) .
\tag{22}
\]

\emph{5 bis) Les fonctions $\wp - e_i$} (pages 83 et 84).\footnote{La page
numérote cette section « 5) » une seconde fois.} Hermite appliqué à $1/(\wp u -
e_i)$, avec $e_1 + e_2 + e_3 = 0$, donne
\[
\frac{1}{\wp u - e_i} = \frac{\wp(u + \omega_i) - e_i}{(e_j - e_i)(e_k - e_i)}
\qquad (\{i, j, k\} = \{1, 2, 3\}), \tag{23}
\]
et (18) avec $v = \omega_i$ :
\[
\wp u - e_i = -\frac{\sigma(u + \omega_i)\sigma(u - \omega_i)}{\sigma^2u\,
\sigma^2\omega_i} = \Bigl(\frac{\sigma_i u}{\sigma u}\Bigr)^2, \qquad
\sigma_i u = \frac{e^{-\eta_i u}\sigma(u + \omega_i)}{\sigma\omega_i} .
\tag{24--25}
\]
La racine carrée $\sqrt{\wp u - e_i} = \sigma_i u/\sigma u$ est donc méromorphe ;
$\sigma_i$ est paire, $\sigma_i(0) = 1$, et
\[
\sigma_i(u + 2\omega_i) = -\sigma_i(u)\,e^{2\eta_i(u + \omega_i)}, \qquad
\sigma_i(u + 2\omega_j) = \sigma_i(u)\,e^{2\eta_j(u + \omega_j)} \quad (j \neq i),
\tag{26}
\]
si bien que $\sigma_i/\sigma$ a la période $2\omega_i$ et change de signe par $u
\mapsto u + 2\omega_j$ : ses périodes sont engendrées par $2\omega_i$ et
$4\omega_j$ --- pour $\sqrt{\wp - e_1}$, $2\omega_1$ et $4\omega_2$ ; pour
$\sqrt{\wp - e_2}$, $2\omega_2$ et $4\omega_3$ ; pour $\sqrt{\wp - e_3}$,
$2\omega_3$ et $4\omega_1$.

\pagerange{84}{89}
\subsection*{Les fonctions $\Theta$ (pages 84 à 89)}

Elles servent au calcul numérique
et rompent la symétrie entre $\omega_1$ et $\omega_2$. Par $u \mapsto u +
2\omega_1$, $\sigma(u)$ et $e^{\eta_1u^2/2\omega_1 - i\pi u/2\omega_1}$ sont
multipliés par le même facteur $e^{2\eta_1(u + \omega_1) - i\pi}$ ; leur quotient
a la période $2\omega_1$ et se développe en série de Fourier :
\[
\sigma(u) = e^{\eta_1u^2/2\omega_1}\sum_{n \in \mathbf{Z}} A_n\,e^{(2n-1)i\pi
u/2\omega_1},
\]
et la quasi-périodicité en $2\omega_2$ donne, avec la relation de Legendre, la
récurrence $A_{n+1} = -h^{2n}A_n$, $h = e^{i\pi\omega_2/\omega_1}$,
$\Im(\omega_2/\omega_1) > 0$.\footnote{L'exposant $2n$ est lu sous réserve ; le
calcul le confirme.} D'où la fonction
\[
\Theta(v \mid t) = -\frac1i\sum_{n \in \mathbf{Z}}(-1)^n e^{(n - \frac12)^2
i\pi t + (2n-1)i\pi v} = 2\sum_{n \geqslant 1}(-1)^{n-1}h^{(n - \frac12)^2}\sin(2n-1)\pi v,
\qquad h = e^{i\pi t}, \tag{1--2, 6}
\]
holomorphe pour $\Im t > 0$ --- c'est la fonction $\vartheta_1(v \mid t)$ de
Jacobi --- et
\[
\sigma(u \mid \omega_1, \omega_2) = \frac{2\omega_1}{\Theta_v'(0 \mid
\omega_2/\omega_1)}\,e^{\eta_1u^2/2\omega_1}\,\Theta\Bigl(\frac{u}{2\omega_1}
\Bigm| \frac{\omega_2}{\omega_1}\Bigr) . \tag{3}
\]
$\Theta$ est impaire en $v$, ses zéros sont les $v = m_1 + m_2t$, et
\[
\Theta(v + 1 \mid t) = -\Theta(v \mid t), \qquad \Theta(v + t \mid t) =
\Theta(v \mid t)\,e^{-2i\pi v - i\pi t - i\pi}, \qquad \Theta(v \mid t + 4) =
-\Theta(v \mid t) ; \tag{4--5}
\]
les deux premières relations caractérisent $\Theta$ à un facteur près. Pour $t
= iT$, $\Re T > 0$, la série devient $2\sum(-1)^{n-1}e^{-(n - \frac12)^2\pi
T}\sin(2n-1)\pi v$, et pour $\Re a > 0$
\[
\sum_{n \in \mathbf{Z}} e^{-an^2 + bn + c} = -e^{c + \frac b2 - \frac
a4}\,\Theta\Bigl(\frac{b - a - i\pi}{2i\pi} \Bigm| \frac{ia}{\pi}\Bigr) . \tag{7}
\]

\emph{Transformation $t \mapsto -1/t$.} Écrire que $\sigma(u \mid \omega_1,
\omega_2) = \sigma(u \mid \omega_2, -\omega_1)$ --- c'est le même réseau --- et
utiliser la relation de Legendre donne
\[
\Theta(v \mid t) = e^{-i\pi v^2/t}\,\Theta\Bigl(\frac vt \Bigm| -\frac1t\Bigr)\,
\cdot t\,H(t), \qquad H(t) = \frac{\Theta_v'(0 \mid t)}{\Theta_v'(0 \mid -1/t)} .
\tag{8}
\]
Pour expliciter le facteur, on note que $\Theta$ est solution de l'équation de
la chaleur
\[
\frac{\partial^2\Theta}{\partial v^2} = 4i\pi\,\frac{\partial\Theta}{\partial
t} ; \tag{9}
\]
l'appliquer aux deux membres de (8) montre que $tH(t)$ est proportionnel à
$1/\sqrt t$, et le cas $t = i$ fixe la constante :
\[
\boxed{\;\Theta(v \mid t) = e^{-i\pi v^2/t}\,\Theta\Bigl(\frac vt \Bigm|
-\frac1t\Bigr)\, i\sqrt{\frac it}\;}, \qquad
\Theta_v'(0 \mid t) = \frac it\sqrt{\frac it}\;\Theta_v'\Bigl(0 \Bigm|
-\frac1t\Bigr), \tag{10--11}
\]
la racine carrée étant la détermination principale ($\Re\sqrt{i/t} >
0$).\footnote{La page écrit (8) sans le facteur $t$ : la dérivée en $v = 0$ du
second membre fait apparaître $\frac1t\Theta_v'(0 \mid -1/t)$, et le facteur
est nécessaire --- on l'a vérifié numériquement. Elle trouve ensuite
« $H'/H = -1/2t$, d'où $H$ proportionnel à $1/\sqrt t$ », ce qui est vrai de
$tH$ ; et (11), qui donne $H = \frac it\sqrt{\frac it}$, proportionnel à
$t^{-3/2}$, est juste, comme (10) et (12). La constante du cas $t = i$ est
écrite « $a = \sqrt i$ » : pour que (10) en résulte, elle vaut $i\sqrt i$. Les
bornes de l'argument se lisent mal. Le signe de la page entre $\Theta(v/t \mid
-1/t)$ et le quotient ressemble à un $+$ ; c'est un produit.} Pour $t = iT$,
$\Re T > 0$, (11) s'écrit, avec $\Theta_v'(v \mid t) = 2\pi\sum(-1)^{n-1}(2n-1)
e^{(n - \frac12)^2 i\pi t}\cos(2n-1)\pi v$,
\[
\sum_{n \geqslant 1}(-1)^{n-1}(2n-1)e^{-(n-\frac12)^2\pi T} =
\frac{1}{T^{3/2}}\sum_{n \geqslant 1}(-1)^{n-1}(2n-1)e^{-(n-\frac12)^2\pi/T}. \tag{12}
\]

La formule (10) ramène le calcul de $\Theta(v \mid iT)$ à celui de $\Theta(v \mid
i/T)$ ; pour $T$ réel on peut donc supposer $T \geqslant 1$, d'où $|h| \leqslant
e^{-\pi} \approx 0{,}043$, et la série, dont l'exposant de $h$ croît comme $n^2$,
converge très vite.\footnote{La page écrit « $e^{-\pi} \mathbin{\#} \frac{1}{10}$
» (son signe d'approximation) ; $e^{-\pi}$ vaut environ $1/23$, ce qui ne fait
que renforcer l'argument. Pour $T$ complexe, l'alternative « $\Re T > 1$ ou
$\Re\frac1T > 1$ » n'est pas toujours vraie ; il faut alors tout le groupe
modulaire pour ramener $t$ dans le domaine fondamental.}

\emph{Les autres fonctions thêta.} On définit $\Theta_1, \Theta_2, \Theta_3$ par
\[
\sigma_i(u) = e^{\eta_1u^2/2\omega_1}\,\frac{\Theta_i(u/2\omega_1 \mid
\omega_2/\omega_1)}{\Theta_i(0 \mid \omega_2/\omega_1)}, \tag{13}
\]
et l'on trouve
\[
\Theta_1(v \mid t) = 2\sum_{n \geqslant 1}h^{(n-\frac12)^2}\cos(2n-1)\pi v =
\vartheta_2(v),
\]
\[
\Theta_2(v \mid t) = 1 + 2\sum_{n \geqslant 1}(-1)^nh^{n^2}\cos 2n\pi v =
\vartheta_0(v),
\]
\[
\Theta_3(v \mid t) = 1 + 2\sum_{n \geqslant 1}h^{n^2}\cos 2n\pi v =
\vartheta_3(v), \qquad \Theta = \vartheta_1 , \tag{14}
\]
dans les notations de Jacobi ($\vartheta_0$ est le $\vartheta_4$
d'aujourd'hui).\footnote{Vérifié numériquement, avec $\omega_3 = -\omega_1 -
\omega_2$. La page écrit la forme en cosinus de $\Theta_2$ sans le
$(-1)^n$, que sa forme exponentielle, juste, comporte. Elle appelle les
$\vartheta_i$ « fonctions de Weierstrass » ; la notation est de Jacobi.} Une
fois fixés $\omega_1$ et $t = \omega_2/\omega_1$, toutes les quantités de la
théorie se calculent par les thêta :
\[
e_i = \Bigl(\frac{1}{2\omega_1}\Bigr)^2\Bigl(\frac13\frac{\Theta'''(0)}{\Theta'(0)} -
\frac{\Theta_i''(0)}{\Theta_i(0)}\Bigr),
\]
\[
e_1 = \frac13\Bigl(\frac{\pi}{2\omega_1}\Bigr)^2\bigl[\Theta_2(0)^4 +
\Theta_3(0)^4\bigr],
\]
\[
e_2 = -\frac13\Bigl(\frac{\pi}{2\omega_1}\Bigr)^2\bigl[\Theta_1(0)^4 +
\Theta_3(0)^4\bigr],
\]
\[
e_3 = \frac13\Bigl(\frac{\pi}{2\omega_1}\Bigr)^2\bigl[\Theta_1(0)^4 -
\Theta_2(0)^4\bigr], \tag{15}
\]
\[
2\eta_1\omega_1 = -\frac{\Theta'''(0)}{6\,\Theta'(0)}, \qquad \Theta'(0) =
\pi\,\Theta_1(0)\Theta_2(0)\Theta_3(0),
\]
les dérivées étant prises en $v$.\footnote{La page écrit $\frac12$ devant $e_3$ ;
c'est $\frac13$, et le second indice, illisible, est $2$ --- vérifié, comme les
deux autres et les deux relations suivantes.}

\pagerange{90}{90}
\subsection*{Récapitulation (page 90)}

Les données $\omega_1, \omega_2$ (avec
$\omega_1/\omega_2 \notin \mathbf{R}$), $g_2, g_3$ (avec $g_2^3 - 27g_3^2 \neq
0$), $c_1, c_2$, ou $e_1, e_2, e_3$ distincts, se déterminent les unes les
autres :
\[
c_1 = 3\sum_{w \in \mathfrak{P}\smallsetminus 0}\frac{1}{w^4}, \quad c_2 =
5\sum_{w \in \mathfrak{P}\smallsetminus 0}\frac{1}{w^6}, \quad g_2 = 20c_1,
\quad g_3 = 28c_2,
\]
\[
e_i = \wp(\omega_i) = \frac{1}{\omega_i^2} + \sum_{w
\neq 0}\Bigl[\frac{1}{(\omega_i - w)^2} - \frac{1}{w^2}\Bigr] ;
\]
$e_1, e_2, e_3$ sont les racines de $4z^3 - g_2z - g_3 = 4(z - e_1)(z - e_2)(z -
e_3)$, d'où
\[
g_2 = -4(e_2e_3 + e_3e_1 + e_1e_2) = 2(e_1^2 + e_2^2 + e_3^2), \qquad g_3 =
4e_1e_2e_3 ;
\]
et $\omega_i = \int_\infty^{e_i} dz/\sqrt{4z^3 - g_2z - g_3}$, ou $\omega_i =
\int_{e_j}^{e_k}dz/\sqrt{4z^3 - g_2z - g_3}$, pour des déterminations convenables
avec $\omega_1 + \omega_2 + \omega_3 = 0$.\footnote{La page écrit $g_2 = 4(e_2e_3
+ e_3e_1 + e_1e_2)$, sans le signe $-$ ; « $g_? = 20c_?$ » dans la dernière
ligne, pour $g_3 = 28c_2$ ; les sommes des $c_i$ portent sur les $(1/2\omega)$
pour $\omega$ dans un « $\mathfrak{P}^*$ » qui est ici, semble-t-il, l'ensemble
des demi-périodes non nulles, et la série de $e_i$ s'écrit sans les termes
$-1/w^2$ qui la font converger.}

\pagerange{91}{92}
\subsection*{La fonction modulaire (pages 91 et 92)}

Comme fonctions de $(\omega_1,
\omega_2)$, avec $\Im(\omega_2/\omega_1) > 0$, $g_2$ et $g_3$ sont homogènes de
degrés $-4$ et $-6$ ; le quotient
\[
J(\tau) = \frac{g_2^3}{\Delta} = \frac{g_2^3}{g_2^3 - 27g_3^2} =
\frac{1}{1 - 27g_3^2/g_2^3}, \qquad \tau = \frac{\omega_2}{\omega_1},
\]
est homogène de degré $0$, donc une fonction de $\tau$ seul, holomorphe sur le
demi-plan $\Im\tau > 0$ : c'est la fonction modulaire.\footnote{La page dit
« homogènes de degré $2$ et $3$ » : ce sont les degrés en $(g_2, g_3)$ comme
poids, et les degrés d'homogénéité en $(\omega_1, \omega_2)$ sont $-4$ et
$-6$ ; le quotient est de degré $0$ dans les deux cas.} Comme $g_2, g_3$ peuvent
être pris arbitraires sous la condition $\Delta \neq 0$, $J$ prend toutes les
valeurs complexes ; comme $g_2, g_3$ ne dépendent que du réseau, $J$ est
invariante par le groupe modulaire des $\tau \mapsto (a\tau + b)/(c\tau + d)$,
$a, b, c, d \in \mathbf{Z}$, $ad - bc = 1$ ; et $J(\tau) = a$ détermine le réseau
à homothétie près, donc $\tau$ à une transformation modulaire près. Il en
résulte que $J$ est biunivoque au voisinage de tout point qui n'est point fixe
d'aucune transformation modulaire non triviale, et en ces points seulement ; il
y a exactement deux classes de tels points, $i$ et $\rho = \frac12 +
i\frac{\sqrt3}{2}$, et
\[
J(i) = 1 \quad (g_3 = 0), \qquad J(\rho) = 0 \quad (g_2 = 0) .
\]
Ce sont les seuls points singuliers de la fonction inverse $\tau(J)$, et ces
propriétés sont à la base de la démonstration des théorèmes de Picard sur les
fonctions entières.\footnote{La page apparie $J(i)$ à $0$ et à $g_2 = 0$, et
$J(\rho)$ à $1$ et à $g_3 = 0$ ; c'est l'inverse, et on l'a vérifié (pour le réseau
carré, $\wp$ est « lemniscatique », $g_3 = 0$). Les valeurs $0$ en $i$ et $1$ en
$\rho$ sont celles de $1 - J$, c'est-à-dire de l'invariant $j$ des pages 39 à
47 : pour celui-là, $j = 0$ est le cas de $\mathbf{Z}/2$ (page 47) et $j = 1$ celui de
$\mathbf{Z}/3$. L'écart vient peut-être de là ; la page ne le dit pas.}

\pagerange{93}{96}
\subsection*{Les fonctions elliptiques de Legendre (pages 93 à 96)}

$\mathrm{sn}\,u$
est la fonction réciproque de $\int_0^u dx/\sqrt{(1-x^2)(1-k^2x^2)}$, impaire ;
on pose, avec $k'^2 = 1 - k^2$,
\[
\mathrm{cn}^2 + \mathrm{sn}^2 = 1, \quad \mathrm{dn}^2 + k^2\mathrm{sn}^2 = 1,
\quad \mathrm{dn}^2 - k^2\mathrm{cn}^2 = k'^2, \qquad \mathrm{sn}\,0 = 0,\
\mathrm{cn}\,0 = \mathrm{dn}\,0 = 1, \tag{1}
\]
$\mathrm{cn}$ et $\mathrm{dn}$ étant uniformes et paires ; les équations
différentielles (2) sont celles de la page 69, et les développements de Taylor
en $0$ sont, avec $\alpha = \frac12(k + \frac1k)$\footnote{Coefficients vérifiés. En termes de $k$, $2\alpha k = 1 + k^2$,
$4k^2(\alpha^2 + 3) = 1 + 14k^2 + k^4$, $8k^3(\alpha^3 + 33\alpha) = 1 + 135k^2 +
135k^4 + k^6$.},
\[
\mathrm{sn}\,u = u - \frac{2\alpha k}{3!}u^3 + \frac{4k^2(\alpha^2 +
3)}{5!}u^5 - \frac{8k^3(\alpha^3 + 33\alpha)}{7!}u^7 + \cdots,
\]
\[
\mathrm{cn}\,u = 1 - \frac{u^2}{2!} + \frac{1 + 4k^2}{4!}u^4 - \frac{1 + 44k^2 +
16k^4}{6!}u^6 + \cdots,
\]
\[
\mathrm{dn}\,u = 1 - \frac{k^2u^2}{2!} + \frac{k^2(4 + k^2)}{4!}u^4 -
\frac{k^2(16 + 44k^2 + k^4)}{6!}u^6 + \cdots . \tag{3}
\]
Pour $k = 0$, $\mathrm{sn} = \sin$, $\mathrm{cn} = \cos$,
$\mathrm{dn} = 1$ ; pour $k = 1$, $\mathrm{sn} = \mathrm{th}$, $\mathrm{cn} =
\mathrm{dn} = 1/\mathrm{ch}$.\footnote{La page écrit $\mathrm{dn} =
1/\mathrm{th}$ pour $k = 1$ ; c'est $1/\mathrm{ch}$, puisque $\mathrm{dn}^2 =
1 - \mathrm{sn}^2 = \mathrm{cn}^2$ quand $k = 1$.}

\emph{Périodes et valeurs remarquables.} Soient
\[
K = \int_0^1\frac{dx}{\sqrt{(1-x^2)(1-k^2x^2)}}, \qquad K' =
\int_0^1\frac{dx}{\sqrt{(1-x^2)(1-k'^2x^2)}}, \qquad K'' = -K - iK' .
\]
\begin{itemize}
\item $\mathrm{sn}$ : périodes $4K$, $2iK'$ ; zéros $2mK + 2m'iK'$ ; pôles $iK' +
2mK + 2m'iK'$ ;
\item $\mathrm{cn}$ : périodes $4K$, $-2K'' = 2K + 2iK'$ ; zéros $K + 2mK +
2m'iK'$ ; pôles $iK' + 2mK + 2m'iK'$ ;
\item $\mathrm{dn}$ : périodes $2K$, $4iK'$ ; zéros $K'' + 2mK + 2m'iK'$ ; pôles
$iK' + 2mK + 2m'iK'$.
\end{itemize}
Les trois fonctions sont d'ordre $2$ pour leurs propres périodes et d'ordre $4$
pour les périodes communes $4K$, $4iK'$. Pour $0 < k < 1$ et $u$ réel,
$\mathrm{sn}$ n'a pas de pôle et s'annule en $2mK$, de période réelle $4K$ ;
$\mathrm{cn}$ n'a pas de pôle, s'annule en $K + 2mK$, de période $4K$ ;
$\mathrm{dn}$ n'a ni pôle ni zéro, de période $2K$.\footnote{Dans le tableau de
la page, les zéros de $\mathrm{cn}$ et de $\mathrm{dn}$ se terminent par
« $2m'iK$ » : c'est $2m'iK'$ ; le mot « pôles » manque devant la dernière ligne
de $\mathrm{dn}$. La note marginale dit de $\mathrm{sn}$ qu'elle « n'a pas de
zéros ni de pôles pour $u$ réel » ; elle a les zéros réels $2mK$.} Les
translations par les demi-périodes\footnote{Toutes vérifiées numériquement. La page donne $\mathrm{dn}(u+K'') =
ik'\,\mathrm{sn}\,u/\mathrm{cn}\,u$ ; le signe est $-$. Son numérateur est
d'ailleurs surchargé.} :
\[
\begin{array}{lll}
\mathrm{sn}(u+2K) = -\mathrm{sn}\,u & \mathrm{cn}(u+2K) = -\mathrm{cn}\,u &
\mathrm{dn}(u+2K) = \mathrm{dn}\,u \\
\mathrm{sn}(u+2iK') = \mathrm{sn}\,u & \mathrm{cn}(u+2iK') = -\mathrm{cn}\,u &
\mathrm{dn}(u+2iK') = -\mathrm{dn}\,u \\
\mathrm{sn}(u+2K'') = -\mathrm{sn}\,u & \mathrm{cn}(u+2K'') = \mathrm{cn}\,u &
\mathrm{dn}(u+2K'') = -\mathrm{dn}\,u \\[1ex]
\mathrm{sn}(u+K) = \dfrac{\mathrm{cn}\,u}{\mathrm{dn}\,u} & \mathrm{cn}(u+K) =
-k'\dfrac{\mathrm{sn}\,u}{\mathrm{dn}\,u} & \mathrm{dn}(u+K) =
\dfrac{k'}{\mathrm{dn}\,u} \\[2ex]
\mathrm{sn}(u+iK') = \dfrac{1}{k\,\mathrm{sn}\,u} & \mathrm{cn}(u+iK') =
-\dfrac ik\dfrac{\mathrm{dn}\,u}{\mathrm{sn}\,u} & \mathrm{dn}(u+iK') =
-i\dfrac{\mathrm{cn}\,u}{\mathrm{sn}\,u} \\[2ex]
\mathrm{sn}(u+K'') = -\dfrac1k\dfrac{\mathrm{dn}\,u}{\mathrm{cn}\,u} &
\mathrm{cn}(u+K'') = -\dfrac{ik'}{k\,\mathrm{cn}\,u} & \mathrm{dn}(u+K'') =
-ik'\dfrac{\mathrm{sn}\,u}{\mathrm{cn}\,u}
\end{array}
\]
La transformation imaginaire de Jacobi :
\[
\mathrm{sn}(iu \mid k) = i\,\frac{\mathrm{sn}(u \mid k')}{\mathrm{cn}(u \mid
k')}, \qquad \mathrm{cn}(iu \mid k) = \frac{1}{\mathrm{cn}(u \mid k')}, \qquad
\mathrm{dn}(iu \mid k) = \frac{\mathrm{dn}(u \mid k')}{\mathrm{cn}(u \mid k')} ;
\]
et la table des valeurs :\footnote{Table vérifiée, case par case, telle qu'elle est sur la page.}
\[
\begin{array}{c|cccccccccc}
 & 0 & K & 2K & 3K & iK' & 2iK' & 3iK' & K'' & 2K'' & 3K''\\
\hline
\mathrm{sn} & 0 & 1 & 0 & -1 & \infty & 0 & \infty & -\frac1k & 0 & \frac1k\\
\mathrm{cn} & 1 & 0 & -1 & 0 & \infty & -1 & \infty & -i\frac{k'}{k} & 1 & -i\frac{k'}{k}\\
\mathrm{dn} & 1 & k' & 1 & k' & \infty & -1 & \infty & 0 & -1 & 0
\end{array}
\]


\emph{Relations avec les fonctions de Weierstrass} (page 95). Supposons
$\omega_1$ réel et $\omega_2$ purement imaginaire, $\Im\omega_2 > 0$ ; les $e_i$
sont réels, et $e_2 < e_3 < e_1$. Posons $k = \sqrt{(e_3 - e_2)/(e_1 - e_2)}$ et
$M = \sqrt{e_1 - e_2}$. Alors
\[
\mathrm{sn}(u \mid k) = M\,\frac{\sigma(u/M)}{\sigma_2(u/M)}, \qquad
\mathrm{cn}(u \mid k) = \frac{\sigma_1(u/M)}{\sigma_2(u/M)}, \qquad
\mathrm{dn}(u \mid k) = \frac{\sigma_3(u/M)}{\sigma_2(u/M)},
\]
\[
\wp(u) = e_2 + \frac{M^2}{\mathrm{sn}^2(Mu \mid k)}, \qquad K = M\omega_1,\quad iK'
= M\omega_2,\quad K'' = M\omega_3 .
\]
Inversement, partant de $\mathrm{sn}(u \mid k)$, on prend $M$ arbitraire, on
définit les $\omega_i$ par ces relations, et
\[
3e_1 = (1 + k'^2)M^2, \qquad 3e_2 = (k'^2 - 2)M^2, \qquad 3e_3 = (1 -
2k'^2)M^2 .
\]
On prend souvent $M = \frac12$, d'où $2\omega_1 = 4K$, $2\omega_2 = 4iK'$, pour
lesquelles $\mathrm{sn}$, $\mathrm{cn}$, $\mathrm{dn}$ sont d'ordre
$4$.\footnote{Vérifié numériquement. L'indice du dénominateur de $\mathrm{dn}$
se lit mal sur la page ; c'est $2$. Le premier terme de $3e_2$ est surchargé (un
$2$ sur un autre signe) : c'est $-2$, puisque $e_1 + e_2 + e_3 = 0$.}

\emph{Fonctions de Legendre et fonctions $\Theta$} (page 96). Avec $v = u/2K$ et
$\tau = iK'/K$\footnote{Les trois expressions par les thêta sont vérifiées numériquement.
Ici la variable de $\Theta$ est $u/2K$, si bien que $u \mapsto u + K$ est $v
\mapsto v + \frac12$.},
\[
\mathrm{sn}(u \mid k) = \frac{1}{\sqrt k}\,\frac{\Theta(v \mid
\tau)}{\Theta_2(v \mid \tau)}, \qquad \mathrm{cn}(u \mid k) =
\sqrt{\frac{k'}{k}}\,\frac{\Theta_1(v \mid \tau)}{\Theta_2(v \mid \tau)},
\qquad \mathrm{dn}(u \mid k) = \sqrt{k'}\,\frac{\Theta_3(v \mid
\tau)}{\Theta_2(v \mid \tau)} ;
\]
et les formules d'addition, sur lesquelles le dossier s'achève :
\[
\mathrm{sn}(u+v) = \frac{\mathrm{sn}\,u\,\mathrm{cn}\,v\,\mathrm{dn}\,v +
\mathrm{sn}\,v\,\mathrm{cn}\,u\,\mathrm{dn}\,u}{1 -
k^2\mathrm{sn}^2u\,\mathrm{sn}^2v}, \qquad
\mathrm{cn}(u+v) = \frac{\mathrm{cn}\,u\,\mathrm{cn}\,v -
\mathrm{sn}\,u\,\mathrm{sn}\,v\,\mathrm{dn}\,u\,\mathrm{dn}\,v}{1 -
k^2\mathrm{sn}^2u\,\mathrm{sn}^2v},
\]
\[
\mathrm{dn}(u+v) = \frac{\mathrm{dn}\,u\,\mathrm{dn}\,v -
k^2\mathrm{sn}\,u\,\mathrm{sn}\,v\,\mathrm{cn}\,u\,\mathrm{cn}\,v}{1 -
k^2\mathrm{sn}^2u\,\mathrm{sn}^2v} .
\]


\end{document}
